\documentclass[11pt,oneside]{article}
\usepackage[T1]{fontenc}
\usepackage{lmodern}
\usepackage{amsmath,amssymb,amsthm,mathtools}
\usepackage[margin=1in]{geometry}
\usepackage{microtype,booktabs,longtable,array,enumitem}
\usepackage{xcolor}
\usepackage{xurl}
\usepackage[colorlinks=true,linkcolor=blue!45!black,citecolor=blue!45!black,urlcolor=blue!45!black]{hyperref}
\hypersetup{pdftitle={Mathematics in iPRA through finite observations},
pdfauthor={Research report prepared for Vladimir Reshetnikov},
pdfsubject={Finite observation contracts, constructive analytic branches, and algebraic point bounds}}
\setlength{\emergencystretch}{2em}
\setlength{\parskip}{0.3em}
\setlist{itemsep=0.2em,topsep=0.4em}
\newtheorem{theorem}{Theorem}[section]
\newtheorem{lemma}[theorem]{Lemma}
\newtheorem{proposition}[theorem]{Proposition}
\newtheorem{corollary}[theorem]{Corollary}
\theoremstyle{definition}
\newtheorem{definition}[theorem]{Definition}
\newtheorem{example}[theorem]{Example}
\theoremstyle{remark}
\newtheorem{remark}[theorem]{Remark}
\newcommand{\ipra}{\ensuremath{\mathrm{iPRA}}}
\newcommand{\PRA}{\ensuremath{\mathrm{PRA}}}
\newcommand{\PR}{\ensuremath{\mathrm{PR}}}
\newcommand{\NN}{\mathbb N}
\newcommand{\ZZ}{\mathbb Z}
\newcommand{\QQ}{\mathbb Q}
\newcommand{\CC}{\mathbb C}
\newcommand{\abs}[1]{\lvert #1\rvert}
\newcommand{\norm}[1]{\lVert #1\rVert}
\title{\textbf{Mathematics in iPRA through finite observations}\\[0.6em]
\large A general framework, a constructive simple-branch theorem,\\
and finite algebraic bounds on rational points}
\author{Research report prepared for Vladimir Reshetnikov}
\date{15 November 2026}

\begin{document}
\maketitle
\begin{abstract}
We develop a common method for turning mathematical constructions and their
proofs into first-order intuitionistic primitive recursive arithmetic. The
basic unit is a finite observation together with a primitive recursive demand
for input data and a proof using a bounded invariant. We prove bounded
Booleanization and witness selection, an adaptive finite-demand transfer
theorem, composition of observation contracts, and soundness of a restricted
polynomial certificate language. These results separate the mathematical
work from claims that merely a terminating program, a finite checker, or a
native Lean theorem establishes an iPRA theorem.

The framework yields concrete advances on difficulties recorded in the
ProveIt proof collection. For a polynomial with primitive-recursive complex
coefficient names, a supplied simple root and a finite derivative-separation
certificate produce an ordinary analytic branch, with explicit coefficient
precision schedules, a convergence disk, finite tails, and residual bounds.
A finite budget theorem synchronizes stability windows for several
bounded-variation sequences and avoids several finite exceptional sets.
For an odd-degree equation $y^2=f(x)$, we give a finite polynomial proof
that $A(x)+B(x)y$ has at most $\deg(A^2-fB^2)$ distinct rational zeros.
More generally, an explicit spanning certificate of size $D$ for a polynomial
quotient gives at most $D$ rational solutions. The first bound discharges a
specific geometric contract left open in the Siegel packet for the indicated
hyperelliptic systems.

The archive includes human proofs, compiler-checked Lean definitions and
finite-search lemmas, an explicit plan for the remaining embedded
derivations, and exact rational validation code. General height bounds,
quantitative approximation theorems, and singular Puiseux selection are
not claimed to follow from this framework.
\end{abstract}

\noindent\textbf{Source baseline.} ProveIt commit
\texttt{27a07921176bccbf630084b959700605783ff4aa}.
The mathematical audit is targeted; it is not an independent verification
of all 47 packets in the collection.\par
\noindent\textbf{Formalization status.} The human arguments include explicit
iPRA induction accounting. The new analytic and geometric theorems have
not yet been implemented as complete \texttt{Prf PRAAx} derivations.
\clearpage
\tableofcontents
\clearpage

\section{What this framework changes}

The repository's many proof packets already exhibit substantial common
structure: rational approximations, finite algebra, explicit inequalities,
primitive recursive algorithms, and an audit of the allowed induction.
They also repeatedly encounter the same missing interfaces. A geometric
argument may produce a finiteness assertion without a coordinate bound.
An analytic argument may require an exact zero test on a Cauchy name.
A Lean argument may prove a statement about natural numbers while leaving
the object-language arithmetic derivation unconstructed. The useful next
step is to make these interfaces reusable and then prove new instances.

\subsection{The proposed unit of mathematical development}

An observation request is a finite datum: a precision, a polynomial degree,
a bounded list of points, a finite matrix, a requested range of indices,
or a finite assertion about one of these. A construction should specify:
\begin{enumerate}
\item a fixed primitive recursive representation of its inputs and outputs;
\item the finite observations of the input used for this request;
\item an explicit primitive recursive output procedure;
\item a finite correctness predicate, including all necessary promises;
\item a proof whose loop invariants become quantifier-free arithmetic.
\end{enumerate}
Theorem~\ref{thm:finite-demand-transfer} supplies the general transfer step.
It allows adaptive input demands and global hypotheses, provided the
induction itself is performed at a fixed finite request. The global
interpretation is recovered afterwards by universal introduction.

This permits quite different objects to share one proof discipline. An
exact rational algebraic object uses finite equality tests. A real or
complex name uses finite coherence and error tests. A function uses an
evaluator together with a quantitative continuity contract. A finite
geometric configuration uses incidence and nondegeneracy certificates.
A count bound uses a contradiction for a supplied list of too many
distinct points. The framework does not collapse these representations
into one stronger existence claim.

\subsection{The main mathematical outcomes}

The principal positive results are the following.
\begin{enumerate}
\item \textbf{A direct arithmetic transfer theorem.}
Bounded truth, bounded witness selection, adaptive finite demand, and
composition are proved explicitly. A restricted proof discipline has an
effective syntactic compilation into iPRA. Its primitive contracts must
already be proved; the theorem does not supply missing mathematics.

\item \textbf{A separated simple branch for Cauchy coefficients.}
Theorem~\ref{thm:simple-branch} replaces exact algebraic coefficient
codes by a different sufficient interface: arbitrary fixed PR complex
coefficient names, a supplied root at $X=0$, and an explicit separation
of the derivative from zero. Finite polynomial iteration produces
individually accessible coefficient names and quantitative convergence.
The derivative witness is the exact point where the singular Puiseux
obstruction is excluded.

\item \textbf{Common finite windows without a global modulus.}
Theorem~\ref{thm:common-window} combines finite variation bounds and
finite exceptional-set bounds by addition of budgets. It gives one
bounded search that satisfies all observations on the same window.
It includes nonmonotone sequences. A halting example proves why this
output cannot generally be promoted to a Cauchy modulus.

\item \textbf{A concrete repair of a Diophantine geometric dependency.}
Theorem~\ref{thm:hyperelliptic-fiber} supplies a sharp finite zero-count
lemma on odd-degree hyperelliptic equations. It replaces a divisor-based
argument in a specified part of the Siegel program. Conditional on the
still-needed finite covering theorem, it gives the point-count bound
$ML$ for the explicit evaluation vectors used in
Corollary~\ref{cor:hyperelliptic-cardinality}.

\item \textbf{A general finite algebra interface.}
Theorem~\ref{thm:finite-algebra-zeros} converts explicit polynomial
spanning identities into a bound on the number of rational solutions.
The certificate can be checked by coefficient arithmetic. Its proof
uses separator polynomials and finite row reduction, and does not
assume a generic algebraic-geometry decision procedure.
\end{enumerate}

The elementary ingredients of these results are established mathematics.
The contribution of this report is their integration into a common iPRA
discipline, the explicit arithmetic and precision details, and their
application to the specific unresolved contracts in the repository.
No claim of priority for a new classical analytic or algebraic theorem
is intended.

\subsection{Three distinctions that preserve the intended theorem}

\paragraph{Counting points versus finding all points.}
A proof that every list of $C+1$ distinct solutions is impossible gives
a cardinality bound. It does not tell a bounded search where the last
solution lies. For an arbitrary decidable predicate, even a bound of
one permits a sole solution at an arbitrarily late index. A global
height bound is additional information. The finite algebra and
hyperelliptic results in this report are count results.

\paragraph{Proving a finite certificate sound versus producing one.}
A coefficient identity can be verified and its consequence can be
proved uniformly in a supplied certificate. This does not show that
every input in a proposed theorem admits such a certificate. For a
uniform positive result, either supply the certificate as input or
prove a primitive recursive constructor and its correctness. This
distinction applies to algebraic splitting, exceptional subspace
covers, root isolation, and proof traces alike.

\paragraph{Native Lean truth versus an embedded arithmetic proof.}
Lean is the metatheory in the intended formalization. The target theory
is represented by its syntax and derivation type. A theorem about
\texttt{Nat}, \texttt{Rat}, or \texttt{Complex} in Lean is useful
mathematics, but the required final endpoint is a value of the repository
type \texttt{Prf PRAAx [] formula}. The archive's formalization plan
keeps these obligations separate, including which files were actually
compiled in this work.

\subsection{How to read and use the report}

Section~\ref{sec:finite-demand-framework} gives the common foundations.
Section~\ref{sec:simplebranch} develops the analytic repair and a finite
mesh method. Section~\ref{sec:stability} gives the stability selector.
Section~\ref{sec:repository-repair} audits the hard packets and develops
the geometric and finite algebra repairs. The final sections specify
the exact Lean interfaces, remaining obligations, reproducible checks,
and the boundaries of the claims.

The manuscript is intended to stand on its own. The repository is used
to identify the target foundations and the particular gaps; no unverified
analytic theorem from another packet is a hidden premise of the new
finite polynomial, stability, or zero-count proofs. The conditional
covering corollary deliberately leaves its covering theorem as an
explicit hypothesis.

\section{A finite-demand framework for first-order iPRA}
\label{sec:finite-demand-framework}

The central proposed interface is a theorem about a finite observation, accompanied
by an explicit list of the input observations it uses.  Infinite objects may occur
in the mathematical explanation, but the inductive argument is carried by a
primitive recursive state and a decidable invariant.  This is useful in two
ways.  It exposes exactly where an existence theorem needs a bound or a modulus,
and it prevents the considerable metatheory available in Lean from being
silently added to the object theory.

The transfer results below are proved directly in intuitionistic arithmetic.
Their role is narrower than a general conservation theorem: they compile a
specified class of finite arguments, with their hypotheses visible.  They do
not assert that every constructive proof, every computable procedure, or every
Lean theorem has a translation into \ipra.  The ordinary PR presentation and its
closure under bounded predicates are standard; see, for example,
\cite[Section~2]{Avigad2002}.  All the finite-demand arguments given here are
self-contained and do not invoke that paper's higher-type or classical
conservation results.

\subsection{The target theory and the permitted induction}

For this chapter, \ipra\ means first-order intuitionistic predicate logic with
equality, the natural-number constructors $0,S$, the successor axioms
\[
 Sx\ne0,\qquad Sx=Sy\longrightarrow x=y,
\]
and a function symbol for each fixed primitive recursive construction.  The
defining equations include composition and
\[
 r(0,\boldsymbol x)=g(\boldsymbol x),\qquad
 r(Sn,\boldsymbol x)=h(n,r(n,\boldsymbol x),\boldsymbol x).
\]
The induction schema is restricted to quantifier-free formulas $A$, with
parameters allowed:
\[
 \bigl(A(0,\boldsymbol x)\land
       \forall n(A(n,\boldsymbol x)\to A(Sn,\boldsymbol x))\bigr)
       \longrightarrow \forall n\,A(n,\boldsymbol x).
\]
Only finitely many of the PR symbols occur in any individual derivation.
A fresh name for a specified PR construction is legitimate in this all-PR
presentation.  This does not license totality of arbitrary terminating
programs.  When working in a smaller finite presentation, the introduction
must instead be justified by its established representation or translation
theorem.

There is no restriction saying that every theorem or every hypothesis must be
quantifier-free.  In particular, a statement such as
$\forall k\,C(a,k)$ is a legitimate hypothesis.  The restriction concerns the
formula to which induction is applied.  This distinction is indispensable:
one may instantiate a global regularity hypothesis at finitely many demanded
indices while using only a quantifier-free state invariant.

Integers, rationals, finite tuples, lists, and polynomials will always mean
natural-number codes.  Their operations are totalized on malformed codes; a
separate validity predicate specifies the intended inputs.  For example, a
rational is a pair $(z,d)$ with $z$ an integer code and $d>0$.  Equality is
\[
 (z,d)\simeq(w,e)\quad\Longleftrightarrow\quad ze=wd,
\]
so neither a quotient type nor a normalization algorithm is necessary.  Signs
and order can be tested by integer cross-multiplication with positive
denominators.  All these tests are primitive recursive.  An implementation
should keep validity predicates in its contracts even when its executable
functions return harmless default values on invalid data.

\subsection{Bounded truth is an arithmetic operation}
\label{sec:bounded-truth}

A \emph{bounded formula} in this chapter is built from equality of PR terms,
intuitionistic propositional connectives, and quantifiers of the form
$\forall u<t$ and $\exists u<t$, where $t$ does not contain the bound variable.
The bound may be any fixed PR term.  In particular, it need not be polynomial,
elementary, or practically computable.

\begin{lemma}[Booleanization of bounded formulas]
\label{lem:bounded-booleanization}
For every fixed bounded formula $B(\boldsymbol x)$ there is a PR function
$b_B(\boldsymbol x)\in\{0,1\}$ and an \ipra\ proof of
\[
 B(\boldsymbol x)\longleftrightarrow b_B(\boldsymbol x)=1.
\]
Consequently bounded formulas are decidable in \ipra, and induction on a
bounded formula is a derived rule of quantifier-free induction.
\end{lemma}

\begin{proof}
The elementary PR development of natural-number equality provides a Boolean
equality test.  Its usual correctness proof uses only the successor axioms and
quantifier-free induction; the same applies to truncated subtraction, order,
and the zero test.  Negation, conjunction, and disjunction are then Boolean
complement, product, and maximum.  For implication use the Boolean function
which returns $1$ when the antecedent test is $0$ and otherwise returns the
consequent test.  Its equivalence to intuitionistic implication is valid
because the antecedent is already decidable.  This step would be invalid for
an arbitrary, undecided proposition.

For bounded universal quantification use a prefix product, and for bounded
existential quantification use a prefix maximum.  Here is an explicit way to
justify the former without covert induction on a quantified formula.  Let
$c(i)\in\{0,1\}$ and define
\[
 p(0)=1,\qquad p(t+1)=p(t)c(t).
\]
Quantifier-free induction proves $p(t)\in\{0,1\}$ and $p(t+1)\le p(t)$.
With an arbitrary fixed $i$, another quantifier-free induction proves
\[
 i<t\land p(t)=1\quad\longrightarrow\quad c(i)=1.
\]
For the converse, fix $b$ and assume $c(i)=1$ for every $i<b$.  Induct on the
quantifier-free formula $t\le b\to p(t)=1$.  In the step with $t<b$, the fixed
hypothesis supplies $c(t)=1$.  Thus
\[
 p(b)=1\quad\longleftrightarrow\quad\forall i<b\;c(i)=1.
\]
The existential case can be proved by the explicit bounded search in
Lemma~\ref{lem:bounded-choice} below, or by the same argument with a running
maximum and a stored witness.  A structural induction on the fixed formula
$B$, in the metatheory, now gives its Boolean function and an \ipra\ proof of
its correctness.  Replacing $B(n)$ by $b_B(n)=1$ in the induction axiom proves
the last assertion.
\end{proof}

The adjective ``fixed'' matters.  The lemma supplies a schema indexed by a
finite formula and its fixed PR vocabulary.  It does not give a single PR truth
function for arbitrary formulas containing arbitrary PR program codes.  That
stronger assertion is false, as shown below.

\begin{lemma}[Bounded search and finite choice]
\label{lem:bounded-choice}
Let $R(\boldsymbol x,y)$ be a fixed PR relation.  There is a PR function
$\ell_R(\boldsymbol x,b)\le b$ such that \ipra\ proves
\begin{align*}
 \ell_R(\boldsymbol x,b)<b
   &\longleftrightarrow \exists y<b\;R(\boldsymbol x,y),\\
 \ell_R(\boldsymbol x,b)<b
   &\longrightarrow R(\boldsymbol x,\ell_R(\boldsymbol x,b)),\\
 y<\ell_R(\boldsymbol x,b)
   &\longrightarrow \neg R(\boldsymbol x,y).
\end{align*}
For a fixed PR bound $B(\boldsymbol x,i)$ and relation
$R(\boldsymbol x,i,y)$, a single PR finite-list function chooses all witnesses
below their bounds whenever
\[
 \forall i<n\;\exists y<B(\boldsymbol x,i)\;R(\boldsymbol x,i,y)
\]
holds.  No choice axiom is used.
\end{lemma}

\begin{proof}
Use $b$ as the sentinel value for failure and define
\[
 \ell(\boldsymbol x,0)=0,
\]
\[
 \ell(\boldsymbol x,b+1)=
 \begin{cases}
   \ell(\boldsymbol x,b),&\ell(\boldsymbol x,b)<b,\\
   b,&\ell(\boldsymbol x,b)=b\text{ and }R(\boldsymbol x,b),\\
   b+1,&\text{otherwise}.
 \end{cases}
\]
This is primitive recursion with decidable branches.  The three maintained
facts are that the stored value is at most the current bound; that a value
strictly below the bound is a successful witness; and that no earlier position
succeeds.  The last fact may be written as a prefix product of the Boolean
negations of $R$, so all three invariants are quantifier-free.

At the next stage, an already successful value is retained.  If none has yet
been found, the new position $b$ is tested.  A successful test stores $b$ and
uses the earlier failure product; an unsuccessful test extends that product.
This proves the invariants by quantifier-free induction.  In particular, the
sentinel occurs exactly when every position below $b$ fails.  Conversely, a
witness below $b$ precludes the sentinel: apply the earlier-failure invariant
to that witness.  This proves the displayed equivalence, constructively.

For finite choice, set
$w_i=\ell_{R(\boldsymbol x,i,-)}(B(\boldsymbol x,i))$ and build the finite
list $(w_0,\ldots,w_{n-1})$ by primitive recursion.  An entirely numerical
encoding is available.  Put
\[
 M=\max\bigl(1,\max_{i<n}B(\boldsymbol x,i)\bigr),\qquad
 W=\sum_{i<n}w_i(M+1)^i.
\]
Since $w_i\le M$, the base-$(M+1)$ digits of $W$ recover the $w_i$, and
$W<(M+1)^n$.  Digit extraction, the finite maximum, and the finite sum are all
PR.  Under the displayed hypothesis, each $w_i$ is strictly below its
corresponding bound and satisfies $R$.  The argument verifies a finite table;
it never chooses an unbounded sequence of witnesses.
\end{proof}

This lemma precisely identifies one frequent repair.  A finite family of
existence assertions causes no difficulty once PR bounds for the witnesses
have been proved.  Merely knowing $\forall i<n\,\exists y\,R(i,y)$ supplies no
such bound, and an unbounded search cannot be relabelled as bounded search.
Likewise, a terminating executable search has not yet been shown to be
primitive recursive: its termination argument must give an appropriate
bound or a PR construction of the output.

\subsection{The local demand theorem}
\label{sec:local-demand-theorem}

We next make the observation interface precise.  Fix the PR functions
\[
 \operatorname{init}(x),\quad
 \operatorname{step}(x,t,s),\quad
 k(x,t,s),\quad q(x,t,s,j),\quad \operatorname{ret}(x,N,s),
\]
and define the run by primitive recursion:
\[
 S(x,0)=\operatorname{init}(x),\qquad
 S(x,t+1)=\operatorname{step}(x,t,S(x,t)).
\]
The numerical parameter $x$ may contain all finite input data and parameters.
A query $q$ is simply a natural-number code for a demanded input fact.  Let
$U(x,q)$ be a fixed bounded predicate describing that fact.  At stage $t$,
the run demands $k(x,t,S(x,t))$ such facts.  Define
\begin{equation}
\label{eq:demand-log}
 L(x,N)\;:\!\Longleftrightarrow\;
 \forall t<N\;\forall j<k(x,t,S(x,t))\;
 U\bigl(x,q(x,t,S(x,t),j)\bigr).
\end{equation}
This is a bounded predicate, hence a single PR Boolean test.  Its query list
may be adaptive: later queries can depend on the already computed state.

\begin{theorem}[Finite-demand transfer]
\label{thm:finite-demand-transfer}
Let $P(x)$, $I(x,N,t,s)$, and $Q(x,N,z)$ be bounded predicates.  Suppose the
following three implications have \ipra\ derivations:
\begin{align}
 P(x)&\longrightarrow I(x,N,0,\operatorname{init}(x)),
   \label{eq:demand-base}\\
 t<N\land I(x,N,t,s)\land
 \bigl(\forall j<k(x,t,s)\;U(x,q(x,t,s,j))\bigr)
 &\longrightarrow
 I(x,N,t+1,\operatorname{step}(x,t,s)),
   \label{eq:demand-step}\\
 I(x,N,N,s)&\longrightarrow Q(x,N,\operatorname{ret}(x,N,s)).
   \label{eq:demand-finish}
\end{align}
Then \ipra\ proves
\begin{equation}
\label{eq:demand-conclusion}
 P(x)\land L(x,N)\quad\longrightarrow\quad
 Q\bigl(x,N,\operatorname{ret}(x,N,S(x,N))\bigr).
\end{equation}
The induction required by this proof is quantifier-free.
\end{theorem}

\begin{proof}
Fix $x,N$ and the two hypotheses in the antecedent of
\eqref{eq:demand-conclusion}.  Use the invariant
\[
 J(t)\;:\!\Longleftrightarrow\;
 t\le N\longrightarrow I(x,N,t,S(x,t)).
\]
By Lemma~\ref{lem:bounded-booleanization}, $J$ has a quantifier-free Boolean
form.  Its base follows from \eqref{eq:demand-base}.  To prove its step,
suppose $t+1\le N$.  Then $t<N$, so $J(t)$ gives the current invariant and
$L(x,N)$ gives exactly the finitely many hypotheses needed in
\eqref{eq:demand-step}.  Substitute the defining equation of $S$ to obtain
$J(t+1)$.  When $t+1>N$, the implication defining $J(t+1)$ is immediate.
Order is decidable, so this case distinction is intuitionistically valid.
Quantifier-free induction now gives $J(N)$; apply
\eqref{eq:demand-finish}.

The global statement $L(x,N)$ was a fixed hypothesis during the induction.
We did not induct on ``all inputs are valid'', on an unbounded universal
assertion, or on the correctness of arbitrary programs.  Every occurrence of
$U$ used by the proof is an entry of the finite demand log.
\end{proof}

\begin{corollary}[Global assumptions can be supplied locally]
\label{cor:global-to-local}
Under the hypotheses of Theorem~\ref{thm:finite-demand-transfer}, \ipra\ proves
\[
 P(x)\land\forall q\,U(x,q)
 \longrightarrow
 \forall N\;Q\bigl(x,N,\operatorname{ret}(x,N,S(x,N))\bigr).
\]
\end{corollary}

\begin{proof}
Fix $N$.  Instantiate the universal input assumption at each query in
\eqref{eq:demand-log}, obtaining $L(x,N)$, and apply the theorem.  Universal
generalization on $N$ completes the argument.  The quantifier-free induction
has already taken place inside the theorem.
\end{proof}

The corollary permits a convenient mathematical language involving globally
regular names.  The theorem is the more informative arithmetic statement:
it remains meaningful for a partial or inconsistent proposed object as long
as the particular finite observations required by this computation are
certified.  This is especially useful for interfaces between independently
developed theorem packets.

\subsection{Prefix contracts and their composition}

Many applications require all observations up to a stage rather than an
arbitrary query list.  Let $V(x,n)$ be a bounded prefix-validity predicate and
suppose \ipra\ proves
\[
 m\le n\land V(x,n)\longrightarrow V(x,m).
\]
An operation with an accuracy parameter $r$ has a \emph{prefix contract} when
there is a PR demand bound $\beta(x,r)$ and an \ipra\ proof of
\[
 V(x,\beta(x,r))\longrightarrow W(x,r,F(x,r)),
\]
where $F$ is PR and $W$ is bounded.  This contract states an actual finite
arithmetic implication.  Calling the output a real number, a function, a
root, a chain, or a geometric object adds no theorem until its intended
observations have been specified by $W$.

\begin{proposition}[Composition of prefix contracts]
\label{prop:prefix-composition}
Suppose an intermediate family of observations $A(x,j)$ satisfies
\[
 V(x,\beta(x,j))\longrightarrow A(x,j),
\]
and a second construction satisfies
\[
 \bigl(\forall j\le\gamma(x,r)\;A(x,j)\bigr)
       \longrightarrow B(x,r).
\]
All displayed matrices and functions are bounded or PR, and both implications
have \ipra\ proofs.  Then \ipra\ proves
\[
 V(x,\delta(x,r))\longrightarrow B(x,r),\qquad
 \delta(x,r)=\max_{j\le\gamma(x,r)}\beta(x,j).
\]
\end{proposition}

\begin{proof}
The finite maximum is PR.  For any $j\le\gamma(x,r)$, its defining property
gives $\beta(x,j)\le\delta(x,r)$.  Monotonicity of prefix validity and the
first contract imply $A(x,j)$.  Bounded universal introduction supplies the
premise of the second contract.  Its conclusion is $B(x,r)$.
\end{proof}

If $\beta(x,j)$ is monotone in $j$, the maximum can be replaced by
$\beta(x,\gamma(x,r))$.  Otherwise the maximum is necessary.  In particular,
a demand transformer must not be composed by blindly substituting precision
parameters when a later proof needs a whole prefix of earlier observations.
Adaptive dependencies are handled by Theorem~\ref{thm:finite-demand-transfer}
instead of by guessing a nonadaptive bound.

\subsection{What a proof compiler may legitimately promise}
\label{sec:proof-compiler}

A small source language can make the preceding discipline systematic.  Its
numeric computations are generated by fixed PR primitives, composition,
decidable conditionals, and primitive recursion on a numerical loop count.
Its assertions are bounded formulas.  Its proof rules are consequence,
sequential composition, decidable case analysis, and the loop rule with a
bounded invariant.  A primitive rule is available only together with an
\ipra\ derivation of its contract.  A mathematical premise such as regularity
of a named function must enter through a stated numerical observation
predicate, not as an unexamined side condition.

\begin{theorem}[Syntactic compilation of the finite proof discipline]
\label{thm:syntactic-compilation}
For every fixed well-formed source derivation built from the preceding rules,
there is an effective construction of an \ipra\ derivation of its concluding
contract.  Every numeric function introduced by the construction comes with
a finite PR definition.  Every inductive proof obligation is discharged by
quantifier-free induction after Booleanization.
\end{theorem}

\begin{proof}
Proceed by induction on the finite source derivation in the metatheory.
A primitive node already contains the required target derivation.  For
sequential composition, substitute the term denoting the first output in the
second contract and use intuitionistic implication elimination.  A
conditional is handled by decidability of its Boolean guard and disjunction
elimination; both branch functions and the combined conditional are PR.
A consequence node is ordinary intuitionistic implication composition.

For a loop, introduce a fresh function symbol with the displayed primitive
recursion equations for its state.  Booleanize the supplied bounded invariant
and apply Theorem~\ref{thm:finite-demand-transfer}.  The local premises used
by the body form the loop's finite demand log.  All list formation, finite
maximum, and Booleanization constructions use composition and primitive
recursion.  Thus the induction constructs an actual finite target derivation
at every node.

This is a construction on syntax.  The source derivation can be checked, its
new PR definition symbols can be emitted, and its target proof can be
assembled without evaluating all programs represented by source syntax.
Nothing in this construction supplies a universal semantic evaluator or a
uniform internal soundness principle for \ipra.
\end{proof}

The theorem is useful as a design specification for an eventual proof
compiler, but it is not a claim that such a compiler has been implemented in
this archive.  In particular, a collection of Lean theorem statements about
the intended semantic functions is not already an implementation of the
syntactic construction.  The output that ultimately matters is a derivation
in the chosen object-theory calculus, or a formally verified translation
into that calculus.

\subsection{Three distinct ways to represent a mathematical family}
\label{sec:three-representations}

The following distinctions remove a subtle source of overly strong claims.

\paragraph{A schema for a fixed PR name.}
Fix a particular PR function $a(k)$ and its definition.  A real-number name,
for example, can be described by rational approximations $a(k)$ and
arithmetical coherence assertions about pairs of indices.  An operation on
that name is another PR definition built from this particular $a$.  The
result is a theorem schema, one instance for every supplied PR definition.
There is no quantifier over all such functions inside first-order \ipra.

\paragraph{A uniformly PR numerical family.}
One may instead fix a jointly PR function $A(e,k)$, where $e$ is an ordinary
natural-number parameter.  Then \ipra\ can prove one theorem uniformly in
$e$.  This genuinely permits quantification over the parameter family, but
it covers only the family given by this one evaluator $A$.  It must not be
identified with an enumeration of all PR functions.

\paragraph{Certified finite observations.}
A third, more general, interface accepts a finite transcript.  For a fixed
register-machine or Turing-machine syntax, the relation
\[
 \operatorname{Trace}(e,u,v,t,c)
\]
asserts that $c$ is a length-at-most-$t$ computation of program $e$ on input
$u$ with output $v$.  Checking it is PR: parse the finite list of
configurations, verify the initial and final configurations, and verify each
adjacent transition by a bounded scan.  Equivalently, bounded execution with
an explicit numerical clock $t$ is PR.  A finite-demand theorem may
therefore be uniform in arbitrary program indices when it is conditional on
the relevant finite traces.

This last interface does not assert that every requested trace exists.  To
turn it into an unconditional theorem for a particular family, one must
supply a PR clock or a PR trace generator, and prove its correctness in the
object theory.  Passing the code of an arbitrary PR clock is insufficient:
the clock itself must be evaluated, which recreates the same issue.

\begin{proposition}[There is no universal PR evaluator]
\label{prop:no-universal-pr-evaluator}
There is no primitive recursive function $E(e,n)$ such that, for every unary
PR function $f$, some index $e$ satisfies
$E(e,n)=f(n)$ for all $n$.
\end{proposition}

\begin{proof}
If $E$ were PR, so would be $d(n)=E(n,n)+1$.  By the assumed universality,
there would be an index $e$ with $E(e,n)=d(n)$ for every $n$.  At $n=e$ this
gives $E(e,e)=E(e,e)+1$, a contradiction.
\end{proof}

This is a metamathematical diagonal argument about actual PR functions.  It
explains why even a perfectly decidable syntax for PR programs does not
justify a single PR evaluation operation for them all.  The distinction is
not a technical inconvenience to be hidden by notation; it is forced by the
mathematics.

Two useful operations nevertheless remain uniform.  A compiler can
transform finite program syntax into finite program syntax, and a checker
can verify a finite computation with a supplied clock.  Neither operation
evaluates arbitrary PR programs to completion.  The proposed framework uses
these two uniform operations together with fixed-name theorem schemas.

\subsection{Presentations, equality, and observations}
\label{sec:presentations}

A \emph{finite arithmetic presentation} consists of a PR validity predicate
$D(a)$, a PR equivalence test $a\simeq b$ on valid codes, and PR operation
codes with object-theory proofs of their required laws.  Rationals,
finite-dimensional rational vectors, finite polynomial expressions, and
finite combinatorial structures fit this format directly.  This does not
mean that every abstract countable structure has such a presentation; the
operations, validity tests, and equality laws are hypotheses to be supplied.

For completion-type objects, exact equality is usually an infinite relation.
A name $a$ can instead have observations $a[r]$ and finite consistency
predicates $V(a,n)$.  The informal assertion that it is a valid completed
object is
\[
 \operatorname{Valid}(a)\;:\!\Longleftrightarrow\;\forall n\;V(a,n).
\]
Here $a[r]$ means either a fixed PR name, a member of a uniformly PR family,
or a requested entry of a finite transcript, as distinguished above.
Arithmetic operations are supplied by finite contracts.  Equality of two
names is an observational assertion, such as an error inequality at every
accuracy; it is never assumed decidable merely because each individual
accuracy test is decidable.

For example, suppose a fixed PR family of rational approximations satisfies
\[
 \forall p,q\quad
 \abs{a[p]-a[q]}\le 2^{-p}+2^{-q}.
\]
For a finite demand $N$, all instances with $p,q\le N$ form a bounded
predicate after denominators are cleared.  A proof that a computed output
has a desired error at accuracy $r$ can request an appropriate finite
prefix.  Once that local contract is proved, the global coherence hypothesis
supplies it by instantiation.  There is no need to induct on the displayed
universal assertion or to construct a quotient of all Cauchy sequences
inside the arithmetic.

The same pattern handles many other mathematical concepts.  A continuous
map needs an approximation procedure and a supplied continuity modulus.
An invertible element needs an explicit inverse or an apartness bound
sufficient for the desired computation.  A compact set used for a search
needs a finite net with an explicit mesh and coverage contract.  An algebraic
extension can be presented by finite arithmetic data, but its field laws
still require a proof; a finite multiplication table over an infinite base
field is not automatically a field certificate.  These are interfaces with
mathematical content, not declarations that erase the original hypotheses.

\subsection{A reusable algebraic application: polynomial certificates}
\label{sec:polynomial-certificates}

The restriction against a universal PR evaluator does \emph{not} prohibit
uniform evaluation of a fixed expression grammar.  This distinction gives a
particularly useful reusable module for algebraic and geometric packets.

Fix a PR-presented commutative ring: its carrier predicate $D$, equality
$\simeq$, and operations $0,1,+,-,\cdot$ are fixed PR data, and their carrier-closure, ring,
and congruence laws have \ipra\ proofs.  Expressions are finite trees built
from variables, integer constants, addition, negation, and multiplication.
Let $\operatorname{Eval}(e,a)$ evaluate the tree using a finite assignment
$a$ of carrier elements.  Let $\operatorname{NF}(e)$ be the usual sparse
commutative polynomial obtained by expanding products, representing each
monomial by a finite exponent vector (padding with zeros when necessary),
ordering these vectors, and collecting equal monomials.  Integer coefficients are stored as integer codes.

\begin{proposition}[Uniform soundness of a fixed algebraic certificate language]
\label{prop:polynomial-normalization}
There are PR functions $\operatorname{Eval}$ and $\operatorname{NF}$, uniform
in expression codes and finite assignments, for which \ipra\ proves
\[
 \operatorname{Good}(e,a)\quad\longrightarrow\quad
 \operatorname{Eval}(e,a)\simeq
 \operatorname{EvalPoly}(\operatorname{NF}(e),a).
\]
Consequently equality of normal forms is a sound PR certificate for equality
of the two expressions, on valid assignments.
\end{proposition}

\begin{proof}
Choose a tree coding in which every immediate subtree has smaller numerical
code than its parent.  Such a coding is obtained, for example, by positive
tagging of a pairing function that dominates its arguments.  Parsing and
checking this property are bounded operations.  Evaluation can now be
computed by course-of-values recursion on the numerical code: keep a finite
table of values for all smaller codes, and use the already stored child
values at the current code.  Invalid codes receive the default value $0$.
The table step uses only the fixed ring operations, so the table and its
lookup operation are PR.  This may be an inefficient evaluation method,
but its PR status is immediate.

Normal forms are computed by the same table construction.  A constant is a
one-term polynomial, and a variable is a one-term monomial.  Addition
concatenates the two coefficient lists and collects equal exponent vectors.
Negation changes coefficient signs.  Multiplication forms the finite list
of pairwise products: coefficients multiply and exponent vectors add.
Insertion sorting and collection are finite recursions, hence PR.  The
lengths and coefficient sizes can grow very rapidly, but all the stated
operations are finite compositions and primitive recursions.

We first verify the finite-list identities used by these operations.
Associativity proves that concatenation evaluates as a sum.  Commutativity
proves that adjacent swaps preserve the sum, so insertion sorting preserves
it.  Distributivity and repeated addition prove that merging equal monomials
adds their integer coefficients without changing evaluation.  For products,
induction on the first list and then on the second list gives
\[
 \left(\sum_{i<m}u_i\right)\left(\sum_{j<n}v_j\right)
       =\sum_{i<m}\sum_{j<n}u_i v_j.
\]
The induction formulas compare two coded ring elements, so they are
quantifier-free.  The identity $a^{r+s}=a^r a^s$ is obtained by
quantifier-free induction on $s$; applying it to the finitely many variables
shows that adding exponent vectors multiplies monomial evaluations.
Negative coefficients are handled by the additive inverse laws.  Each
identity is interpreted using $\simeq$, with the supplied congruence laws
justifying substitution.

Finally, prove the normalization claim simultaneously for all valid
expression codes below a numerical cutoff $b$.  The invariant is
\[
 \forall e<b\quad
 \bigl(\operatorname{Good}(e,a)\to
 \operatorname{Eval}(e,a)\simeq
 \operatorname{EvalPoly}(\operatorname{NF}(e),a)\bigr).
\]
It is a bounded formula and therefore has a quantifier-free Boolean form.
At a new code, every child code is smaller, so the induction hypothesis
applies to all children.  The finite-list identities just proved establish
the claim for the corresponding constructor.  An invalid code imposes no
obligation.  This proves the invariant by permitted induction, and hence the
claim for every expression code.  If two normal forms are equal as finite
codes, substitute that equality and use symmetry and transitivity of
$\simeq$ to obtain equality of the original expressions.
\end{proof}

This is an internal soundness theorem for one restricted certificate
language, justified by a specific PR evaluation and a specific bounded
invariant.  It is not uniform reflection for arithmetic proofs.  The
expression grammar has a fixed finite set of already defined operations;
it contains no constructor that executes an arbitrary coded primitive
recursion program.

The proposition is a common foundation for exact trigonometric polynomial
identities, rational-coordinate geometry identities, determinant identities
at a fixed dimension, and many of the repository's certificate-based
arguments.  It proves \emph{soundness} of certificates.  A uniform theorem
about all dimensions or all degrees additionally needs a PR certificate
generator and an \ipra\ proof that its outputs pass the checker, or a direct
uniform algebraic argument.  Running a generator on finitely many examples
cannot establish that extra statement.

\subsection{The Lean boundary and the object-theory goal}
\label{sec:lean-boundary}

A complete formalization should separate four layers.

\begin{enumerate}
\item \textbf{Executable finite arithmetic.}  Implement natural-number
codes, rational and integer operations, finite lists, Boolean predicates,
state recursions, and demand transformers.  These are ordinary Lean
functions at this stage.
\item \textbf{Mathematical correctness.}  Prove the local numerical
contracts and the finite-demand lemmas in Lean.  This validates the
mathematical arguments, but Lean's unrestricted induction is not itself an
\ipra\ derivation.
\item \textbf{PR syntax and bounded formulas.}  Give every used executable
function an explicit construction in a syntactic PR language; relate that
construction to the executable function.  Define object-language terms,
formulas, substitution, the quantifier-free test, and derivations.  Compile
bounded predicates to Boolean PR terms with derivations of the
corresponding equivalences.
\item \textbf{Object-theory derivations.}  Construct the syntactic proof of
the relevant arithmetic contract.  If the final target is ConciseIPRA,
compose with an established translation into its native derivability
relation, with all hypotheses of that translation supplied.
\end{enumerate}

The principal compiler theorem has the schematic target
\[
 \operatorname{SourceDerivation}(d,\Gamma,\varphi)
 \longrightarrow
 \operatorname{Derivation}_{\ipra}
       (\operatorname{compile}(d),\Gamma^*,\varphi^*),
\]
where the stars denote an explicitly defined syntax translation.  Its
conclusion is a proposition \emph{in Lean about finite derivation objects}.
The step from that conclusion to a native ConciseIPRA derivation must use a
proved bridge, not an assumed semantic reflection rule.

For the finite-demand theorem, the first genuinely arithmetic proof object
should use precisely the Booleanization of
\[
 t\le N\to I(x,N,t,S(x,t))
\]
as its induction formula.  Its base and step are obtained by substitution
from the supplied base and step derivations; its final inference substitutes
$N$ and uses the finishing derivation.  These instructions fix the intended
object-level proof rather than merely asking a theorem prover to rediscover
it.

For bounded search, the proof object should carry the running sentinel,
its success test, and the prefix-failure product.  For prefix composition,
it should carry the PR finite maximum and its elementary upper-bound
lemma.  For polynomial normalization, it should carry the bounded-prefix
invariant over subexpression codes displayed above.  These are concrete
main lemmas with finite syntax, explicit induction formulas, and clear
implementation paths.

Finally, the assertion that a proof checker is PR does not justify
\[
 \operatorname{Check}(p,\ulcorner\varphi\urcorner)=1\longrightarrow\varphi
\]
uniformly for all arithmetic formulas inside the theory being checked.
Such a statement is an additional reflection principle; at falsity it
already supplies the theory's own consistency assertion.  Restricted
reflection for a particular finite algebraic language can be established by
the preceding normalization proof.  General arithmetic reflection cannot
be silently substituted for that proof.  This boundary is essential when
reusing repository results whose Lean statements already live in a rich
ambient metatheory.

\section{A positive repair for Puiseux data: a separated simple branch}
\label{sec:simplebranch}

The Puiseux packet in the repository identifies a precise obstruction:
arbitrary Cauchy presentations of the coefficients do not uniformly supply
individually accessible exact coefficients of a Puiseux root. Its example
\(Y^2-X^2-a\) makes the distinction between \(a=0\) and \(a>0\) visible in
one returned coefficient. Finite exact algebraic presentations avoid that
particular problem, but they are not the only usable input interface.
Here we establish a second interface: a supplied root of the constant
polynomial together with a finite certificate separating its derivative
from zero. The coefficients may otherwise be arbitrary primitive
recursive complex names. The result gives one ordinary power-series
branch, with ramification integer \(1\), explicit coefficient evaluators,
and an explicit disk of convergence. It does not address the singular
case where that derivative vanishes.

The proof has a deliberately finite core. Every estimate first concerns
finite polynomials with Gaussian-rational coefficients. Cauchy names are
introduced only after those estimates have been proved. This order makes
the quantifier-free induction requirements visible and avoids using a
fixed-point theorem or analytic completeness as an arithmetic axiom.
The repository motivation is
\texttt{algebra/puiseux\_ipra/evidence/scope\_obstruction.md} in the
proof collection \cite{ProveIt}.

\subsection{The finite polynomial kernel}

For a Gaussian rational \(z=u+iv\), put
\(\lVert z\rVert_1=|u|+|v|\). This norm is submultiplicative:
expanding the two coordinates of a product and applying the rational
triangle inequality gives
\(\lVert zw\rVert_1\leq\lVert z\rVert_1\lVert w\rVert_1\).
For a finite polynomial \(U(X)=\sum_k u_kX^k\) and a positive rational
\(r\), define its weighted coefficient norm by
\[
  \lVert U\rVert_r=\sum_k\lVert u_k\rVert_1r^k.
\]
Finite convolution gives
\[
 \lVert U+V\rVert_r\leq\lVert U\rVert_r+\lVert V\rVert_r,
 \qquad
 \lVert UV\rVert_r\leq\lVert U\rVert_r\lVert V\rVert_r.
 \tag{A1}\label{eq:branch-norm}
\]
These are inequalities between rational numbers. To prove the product
inequality, expand the coefficient at each degree, use
\(r^{a+b}=r^ar^b\), and rearrange the resulting finite double sum.
Induction on the lengths of the arrays proves each rearrangement.
There is no assertion about an infinite sum in (A1).

Let
\[
 Q(X,Z)=\sum_{i+j\leq D}q_{ij}X^iZ^j,
 \qquad q_{00}=q_{01}=0,
 \qquad D\geq2,
 \qquad \sum_{i,j}\lVert q_{ij}\rVert_1\leq B,
\]
where \(B\geq1\) is rational. Define
\[
 s=\frac{1}{4DB},\qquad
 r=\frac{s}{4B}=\frac{1}{16DB^2}.
 \tag{A2}\label{eq:branch-radii}
\]
In particular, \(0<r\leq s\leq1\). All zero coefficients in this
description are specified positions in a finite array. No procedure is
asked to decide whether a represented coefficient vanishes.

\begin{lemma}[A finite contraction estimate]
\label{lem:finite-poly-contraction}
If \(\lVert U\rVert_r,\lVert V\rVert_r\leq s\), then
\[
 \lVert Q(X,U)-Q(X,V)\rVert_r
       \leq\tfrac14\lVert U-V\rVert_r,
 \qquad
 \lVert Q(X,U)\rVert_r\leq s/2.
 \tag{A3}\label{eq:finite-poly-contraction}
\]
\end{lemma}
\begin{proof}
For every positive integer \(j\), the polynomial identity
\[
 U^j-V^j=(U-V)\sum_{h<j}U^hV^{j-1-h}
\]
follows by cancellation of consecutive terms. Its weighted norm is at
most \(j s^{j-1}\lVert U-V\rVert_r\). In a term of \(Q\) that depends
on \(Z\), the excluded position \((0,1)\) ensures
\(i+j-1\geq1\). Consequently
\[
 r^i s^{j-1}\leq s,
 \quad
 \sum_{i\geq0,\,j\geq1}j\lVert q_{ij}\rVert_1r^i s^{j-1}
 \leq DBs=\tfrac14.
\]
Here the sum is over terms with \(j\geq1\); the index \(i\) may be
zero. This proves the first inequality. The other excluded position,
\((0,0)\), gives
\(\lVert Q(X,0)\rVert_r\leq Br=s/4\). Applying the first inequality
with \(V=0\) gives
\(\lVert Q(X,U)\rVert_r\leq s/4+s/4=s/2\).
\end{proof}

Define the finite polynomial iterates
\[
 U_0=0,\qquad U_{n+1}=Q(X,U_n).
 \tag{A4}\label{eq:finite-iterates}
\]
Their degrees can be bounded by \(D^n\), so they can be constructed by
primitive recursion on finite coefficient arrays. If only coefficients
through degree \(K\) are required, truncating after every operation to
that degree gives exactly the same retained coefficients. Truncation
discards only terms of higher degree, and every exponent in \(Q\) is
nonnegative.

\begin{lemma}[Stabilization by degree]
\label{lem:branch-stabilization}
Every \(U_n\) has zero constant coefficient and
\(\lVert U_n\rVert_r\leq s\). Moreover
\[
  X^{n+1}\mid U_{n+1}-U_n.
 \tag{A5}\label{eq:degree-stabilization}
\]
Thus, if \(c_k=[X^k]U_k\) for \(k\geq1\), the coefficient of
\(X^k\) in every \(U_N\), \(N\geq k\), is \(c_k\).
\end{lemma}
\begin{proof}
The zero constant coefficient and the norm bound follow by induction,
using \(q_{00}=0\), \(U_0=0\), and (A3). Suppose \(U,V\) have zero
constant coefficient and \(X^n\mid U-V\). The factorization used in
the preceding proof shows that
\(Q(X,U)-Q(X,V)=(U-V)R(X)\). Every term of \(R\) has a factor
\(X\): if \(i\geq1\), that factor is explicit; if \(i=0\), then
\(j\geq2\), and each product in
\(\sum_{h<j}U^hV^{j-1-h}\) contains at least one factor having zero
constant coefficient. Hence \(X^{n+1}\mid Q(X,U)-Q(X,V)\).
The initial difference \(U_1-U_0=Q(X,0)\) is divisible by \(X\).
Induction gives (A5). Summing the finitely many later differences
proves the assertion about each retained coefficient.
\end{proof}

Put \(S_K(X)=\sum_{1\leq k\leq K}c_kX^k\), with \(S_0=0\).

\begin{theorem}[Finite majorant, tail, and residual]
\label{thm:finite-branch}
For every \(K\leq N\),
\[
 \sum_{1\leq k\leq N}\lVert c_k\rVert_1r^k\leq s.
 \tag{A6}\label{eq:branch-majorant}
\]
If \(\lVert x\rVert_1\leq r/2\), then
\[
 \begin{split}
 \lVert S_N(x)-S_K(x)\rVert_1&\leq s\,2^{-(K+1)},\\
 \lVert Q(x,S_K(x))-S_K(x)\rVert_1&\leq s\,2^{-K}.
 \end{split}
 \tag{A7}\label{eq:branch-tail}
\]
The polynomial \(Q(X,S_K)-S_K\) has zero coefficients through
degree \(K\).
\end{theorem}
\begin{proof}
By stabilization, \(S_N\) is the truncation of \(U_N\), so (A6)
follows from its norm bound. For the first estimate in (A7), expand the
finite tail. For \(k>K\),
\(\lVert x\rVert_1^k\leq r^k2^{-k}\leq r^k2^{-(K+1)}\).
Summing and using (A6) gives the claimed bound.

The coefficients through degree \(K\) of \(Q(X,S_K)\) agree with
those of \(Q(X,U_K)=U_{K+1}\). Stabilization identifies these with
the coefficients of \(S_K\). Write
\(R_K=Q(X,S_K)-S_K\). Its low coefficients therefore vanish, while
(A3) gives
\(\lVert R_K\rVert_r\leq s/2+s\leq2s\). Every remaining degree
is at least \(K+1\). The same finite-tail estimate gives
\[
 \lVert R_K(x)\rVert_1
 \leq2^{-(K+1)}\lVert R_K\rVert_r
 \leq s2^{-K}.
\]
This proves all claims without introducing an infinite series.
\end{proof}

\subsection{The precision loss is explicit and independent of iteration count}

For the application to names, a perturbation estimate is essential.
Let \(\widetilde Q\) have the same excluded coefficient positions,
degree bound, and majorant \(B\), and suppose
\(\lVert q_{ij}-\widetilde q_{ij}\rVert_1\leq\delta\) at every
position. Set \(m=(D+1)^2\), a convenient upper bound on the number
of positions. For \(\lVert U\rVert_r\leq s\),
\[
 \lVert Q(X,U)-\widetilde Q(X,U)\rVert_r\leq m\delta,
 \tag{A8}\label{eq:coefficient-perturbation}
\]
because \(r,s\leq1\). Applying (A3) to the two sequences of
iterates yields, with \(E_n=\lVert U_n-\widetilde U_n\rVert_r\),
\[
 E_0=0,\qquad E_{n+1}\leq E_n/4+m\delta,
 \qquad E_n\leq2m\delta.
 \tag{A9}\label{eq:iterate-perturbation}
\]
The last inequality is proved by induction: its right-hand side is
preserved because \(2m\delta/4+m\delta\leq2m\delta\).
Thus the dependence on the coefficient input does not get worse as the
number of iterations grows.

Let \(e_p=2^{-p}\). A regular complex name in this subsection is a
primitive recursive Gaussian-rational function \(a(p)\) satisfying
\[
 \lVert a(p)-a(q)\rVert_1\leq e_p+e_q.
 \tag{A10}\label{eq:branch-name}
\]
We work with a fixed primitive recursive family of such functions; this
is a schema, not a universal evaluator for all primitive recursive
programs. Assume each coefficient \(q_{ij}\) is supplied in this form,
with the two excluded entries identically zero. A rational majorant
valid at \emph{every} precision is available directly from the input:
one may take
\[
 B=\max\left(1,\sum_{i+j\leq D}
                 (\lVert q_{ij}(0)\rVert_1+2)\right).
 \tag{A11}\label{eq:uniform-majorant}
\]
Indeed (A10) bounds \(\lVert q_{ij}(p)\rVert_1\) by
\(\lVert q_{ij}(0)\rVert_1+1+e_p\leq\lVert q_{ij}(0)\rVert_1+2\).
A sharper certified majorant can be used instead. This distinction
matters: the finite proof now applies literally to the rational
polynomial \(Q_p\) at each precision, without an informal passage from
exact coefficient inequalities to floating-point ones.

For a nonnegative rational \(u\), write \(\lambda(u)\) for the least
\(h\leq\lceil u\rceil+1\) such that
\(2^h\geq\max(1,u)\). This is a bounded primitive recursive search;
\(2^h\geq h+1\) proves that the indicated bound suffices. Let
\(U_k^p\) be the \(k\)-th iterate for \(Q_p\). An explicit name for
the \(k\)-th coefficient is
\[
 \begin{split}
 h_k&=\lambda(2m r^{-k}),\\
 p(k,n)&=n+h_k+2,\\
 C(k,n)&=[X^k]U_k^{p(k,n)}.
 \end{split}
 \tag{A12}\label{eq:coefficient-name}
\]
For \(k=0\), set \(C(0,n)=0\). The coefficient estimate from
(A9) gives
\[
 \lVert [X^k]U_k^p-[X^k]U_k^q\rVert_1
 \leq2m r^{-k}(e_p+e_q).
\]
Substituting (A12) bounds the difference between \(C(k,n)\) and
\(C(k,n')\) by \((e_n+e_{n'})/4\). Therefore these are regular
names, uniformly in the numerical index \(k\). All loops may be
truncated at degree \(k\).

There is also a uniform evaluator for the finite partial sums. Suppose
\(\lVert x\rVert_1,\lVert x'\rVert_1\leq r/2\), and
\(\lVert U\rVert_r\leq s\). Factoring \(x^k-(x')^k\) gives
\[
 \lVert U(x)-U(x')\rVert_1
 \leq \frac{s}{r}\lVert x-x'\rVert_1.
 \tag{A13}\label{eq:polynomial-eval-lipschitz}
\]
To check the constant, the \(k\)-th summand is bounded by
\(\lVert u_k\rVert_1 k(r/2)^{k-1}\lVert x-x'\rVert_1\),
and \(k2^{-(k-1)}\leq1\) for every \(k\geq1\).

Let a name \(x(p)\) satisfy both (A10) and
\(\lVert x(p)\rVert_1\leq r/2\) at every precision. Put
\[
 T_K(p)=\bigl(\text{truncation of }U_K^p
                    \text{ through degree }K\bigr)(x(p)),
 \qquad C_* =2m+s/r.
\]
Combining (A9), truncation, and (A13) gives the entirely finite estimate
\[
 \lVert T_K(p)-T_K(q)\rVert_1
       \leq C_*(e_p+e_q),
 \tag{A14}\label{eq:partial-sum-name}
\]
with no dependence of \(C_*\) on \(K\).
The bounded-name presentation of the disk is harmless. For a rational
radius \(R>0\), define the rational radial projection by
\(\pi_R(z)=z\) when \(\lVert z\rVert_1\leq R\), and
\(\pi_R(z)=Rz/\lVert z\rVert_1\) otherwise. It has Lipschitz
constant at most \(2\). To see this, assume
\(\lVert z\rVert_1\geq\lVert w\rVert_1\). If both are outside,
subtract the two scaled vectors, keep \(R/\lVert z\rVert_1\) as
the common factor, and bound the remaining scale difference using
\(\lVert z\rVert_1-\lVert w\rVert_1\leq\lVert z-w\rVert_1\).
If only \(z\) is outside, the same calculation uses
\(\lVert z\rVert_1-R\leq\lVert z\rVert_1-\lVert w\rVert_1\).
If both are inside, the claim is immediate. Thus applying \(\pi_R\)
to a regular name at index \(n+1\) gives a regular name in the disk.
It represents the original value whenever that value is in the disk,
because \(\pi_R\) fixes that value and the same Lipschitz estimate
controls approximation errors. Alternatively, one can require the
bounded presentation directly in the input contract.

If \(y_0\) is a further regular name, define
\[
 \begin{split}
 K(n)&=n+3,\\
 p(n)&=n+\lambda(C_*+1)+3,\\
 Y(n)&=y_0(p(n))+T_{K(n)}(p(n)).
 \end{split}
 \tag{A15}\label{eq:root-name}
\]
This is a regular name. For example, if \(n\leq n'\), compare
\(T_{K(n)}(p(n))\) first with \(T_{K(n)}(p(n'))\), using (A14),
and then with \(T_{K(n')}(p(n'))\), using the finite tail in (A7)
for the one rational polynomial \(Q_{p(n')}\). The sum of these
errors and the error in \(y_0\) is at most
\[
 (C_*+1)(e_{p(n)}+e_{p(n')})
          +s2^{-(n+4)}
 \leq\tfrac18(e_n+e_{n'})+\tfrac1{16}e_n
 \leq e_n+e_{n'}.
\]
The case \(n'\leq n\) is symmetric and is decided by natural-number
comparison. This proof constructs the name; it does not extract it from
an assertion that an unspecified sequence converges.

\subsection{Normalization and inversion from one finite separation certificate}

We now explain how the excluded entries of \(Q\) arise without an exact
zero test. Let
\(P(X,Y)=\sum_{i+j\leq D}a_{ij}X^iY^j\), let \(y_0\) be a
supplied regular name, and assume the supplied arithmetic contracts
prove \(P(0,y_0)=0\). Let
\(d=P_Y(0,y_0)\). The additional input is a finite derivative
separation witness. One convenient form consists of \(t\) and a
coordinate \(j\in\{\operatorname{Re},\operatorname{Im}\}\) such that
\[
       |d_j(t)|\geq4e_t.
 \tag{A16}\label{eq:finite-apartness}
\]
This is one decidable rational test, supplied as data; it is not a
conclusion drawn from the negation of \(d=0\).

For completeness, here is an explicit inverse construction. For every
\(p\geq t+2\), regularity gives
\[
 |d_j(p)|\geq4e_t-e_t-e_p\geq\tfrac{11}{4}e_t>2e_t.
\]
Choose a rational \(M\geq1\) bounding
\(\lVert d(t)\rVert_1+2\). Then
\(\lVert d(p)\rVert_1\leq M\) at every precision, and, at the
specified precisions,
\[
 |d(p)|_2^2\geq4e_t^2,
 \qquad
 \lVert d(p)^{-1}\rVert_1\leq
 H:=\frac{M}{4e_t^2}.
\]
The rational inverse is
\((u-iv)/(u^2+v^2)\). Its denominator is strictly positive by the
displayed bound. The identity
\(z^{-1}-w^{-1}=z^{-1}(w-z)w^{-1}\) gives
\[
 \lVert d(p)^{-1}-d(q)^{-1}\rVert_1
             \leq H^2(e_p+e_q).
\]
Consequently
\[
 I(n)=d\bigl(\max(t+2,n+\lambda(H^2)+2)\bigr)^{-1}
 \tag{A17}\label{eq:inverse-name}
\]
is a regular name. The exact rational identity
\(d(p)d(p)^{-1}=1\), together with regularity and the magnitude
bound, proves that it represents an inverse of \(d\). This proof uses
only finite rational identities and explicit error estimates.

Expand the following finite expression:
\[
 Q(X,Z)=Z-d^{-1}P(X,y_0+Z).
 \tag{A18}\label{eq:normalized-branch}
\]
The constant coefficient is zero by the supplied root equation, and
the coefficient at \(Z\) is \(1-d^{-1}d=0\). In the resulting
coefficient array, put these two positions identically equal to the
rational zero. The finite binomial expansion and the supplied equalities
prove that this modified array represents (A18). All other coefficient
names are obtained by finite sums, products, and the inverse (A17).
Here are explicit schedules for those operations. If names \(a,b\)
have rational bounds \(A,B\) at all precisions, addition uses
\(a(n+1)+b(n+1)\), with bound \(A+B\). Multiplication uses
\(a(p)b(p)\), where \(p=n+\lambda(A+B)+1\), with bound \(AB\).
Indeed the product difference is at most
\((A+B)(e_p+e_q)\). A finite expression is evaluated recursively
using these shifts at every node; (A17) handles its only division.
This gives a primitive recursive evaluator and a magnitude bound for
each normalized coefficient. This construction
uses a proof of the two zero identities; it never branches on a test of
equality of arbitrary complex names.

\begin{theorem}[One simple analytic branch with a primitive recursive name]
\label{thm:simple-branch}
For every fixed primitive recursive family of the preceding input data,
the finite certificate (A16), the supplied equality \(P(0,y_0)=0\), and
the regular-name contracts yield the following primitive recursive
outputs and arithmetic contracts.
There is a family of coefficient names \(c_k\), explicitly given by
(A12), with \(c_0=0\), and the rational radii (A2), such that:
\begin{enumerate}
\item For every finite degree \(N\), substitution of
      \(y_0+\sum_{k=1}^{N}c_kX^k\) into \(P\) has zero coefficients
      through degree \(N\).
\item The majorant (A6) holds for every finite coefficient prefix.
      On the disk \(\lVert x\rVert_1\leq r/2\), the finite tail is
      bounded by \(s2^{-(K+1)}\).
\item If \(M\) is a rational bound on \(\lVert d\rVert_1\), then
      \[
       \left\lVert P\left(x,y_0+S_K(x)\right)\right\rVert_1
          \leq Ms2^{-K}.
       \tag{A19}\label{eq:original-residual}
      \]
\item The evaluator (A15) gives a regular root name \(Y(x)\).
      It satisfies \(P(x,Y(x))=0\) in the represented equality
      relation, and \(\lVert Y(x)-y_0\rVert_1\leq s\).
\item It is the unique represented root in that ball.
\end{enumerate}
Thus the output is an ordinary power-series branch; its ramification
integer is \(e=1\).
\end{theorem}
\begin{proof}
The finite polynomial results apply to each rational approximation
\(Q_p\), since the uniform majorant (A11) was chosen for every
precision. The perturbation estimates identify the coefficient names
defined in (A12) and prove that every finite coefficient prefix satisfies
the asserted identities and bounds. In elementary real or complex
notation, this says that \(S_K\) has the stabilized coefficients
established in Lemma~\ref{lem:branch-stabilization}. In arithmetic
notation, each assertion says that its rational approximations differ
by at most the stated error at every requested precision.

The finite low-coefficient identity for \(Q(X,S_N)-S_N\), multiplied
by \(-d\), proves the first claim, using (A18). Multiplication of
the second estimate in (A7) by the magnitude bound \(M\) proves
(A19). Formula (A15) and its finite error calculation give the root
name. To make the last error budget explicit, let \(A_{ij}\) bound
the norms of the original coefficient names at every precision and
let \(A_0\) bound the name of \(y_0\). Put
\[
 R_x=1,\qquad R_y=A_0+s+1,\qquad
 L=\sum_{i+j\leq D,\,j\geq1}
       jA_{ij}R_x^iR_y^{j-1}.
\]
Factoring differences of powers proves
\(\lVert P(x,u)-P(x,v)\rVert_1\leq L\lVert u-v\rVert_1\)
on these bounded disks. The tail and (A19) therefore give
\[
 \lVert P(x,Y(x))\rVert_1\leq s(M+L)2^{-K}
\]
at every finite cutoff \(K\), in the represented inequality
notation. For a requested error \(e_n\), take
\[
 K=n+\lambda\bigl(s(M+L)+1\bigr)+2.
 \tag{A19a}\label{eq:final-zero-budget}
\]
The right-hand side is at most \(e_n/4\). Evaluate the finite
polynomial expressions and their equality certificates with error at
most \(e_n/4\), using the recursive addition, multiplication, and
inverse schedules just specified. This proves the rational zero
certificate with room to spare. All choices are primitive recursive.
The notation of a represented inequality is eliminated by exactly this
requested-precision argument, rather than by an axiom about limits.

Finally, suppose \(z,w\) are two roots after subtracting \(y_0\),
with norms at most \(s\). Use the radial projection described above
to choose equivalent names whose approximants lie in that ball.
The identity (A18) makes them fixed points
of \(Q(x,\cdot)\). The scalar version of (A3) gives
\(\lVert z-w\rVert_1\leq\frac14\lVert z-w\rVert_1\), hence
\(\lVert z-w\rVert_1=0\). At finite precision the same calculation
has an additional known error \(\eta\) and gives
\(\lVert z-w\rVert_1\leq4\eta/3\). Choosing the common input
precision so that this is below \(e_n\) proves the represented
equality. No test of whether either norm is zero occurs.
\end{proof}

\paragraph{What is internal to iPRA.}
The theorem is a schema for fixed evaluator families, with ordinary
numerical parameters for the degree, coefficient indices, precision,
and separation witness. Gaussian rationals, arrays, binomial
coefficients, convolutions, truncations, and finite sums have primitive
recursive codes. The core of (A3)--(A9) is already a collection of
rational inequalities. For a requested final cutoff, all array indices
are bounded by explicitly constructed sizes. Induction on a running
iteration count proves a quantifier-free invariant containing the norm
bound and the bounded list of zero low coefficients. The perturbation
invariant is \(E_n\leq2m\delta\), a rational inequality. Its precision
indices are free numerical parameters during induction. Consequently
no induction formula quantifies over all precisions, all complex
numbers, or all functions. Universally quantified contracts are invoked
only at their finitely many needed numerical instances.

The notation involving represented values in the theorem is an
explanation of these finite contracts. Formalizing those contracts as
\(\mathrm{Prf}\) objects in the repository remains work to be done; a
Lean theorem about complete complex numbers would not establish that
object-language conclusion.

\subsection{Two worked branches and the remaining singular boundary}

For \(P(X,Y)=Y^2-1-X\) and \(y_0=1\), one has \(d=2\) and
\[
 Q(X,Z)=(X-Z^2)/2.
\]
Take \(D=2\), \(B=1\), \(s=1/8\), \(r=1/32\). The explicitly
certified evaluation disk is \(\lVert x\rVert_1\leq1/64\). The
first stabilized coefficients are
\[
 c_1=\frac12,\quad c_2=-\frac18,\quad c_3=\frac1{16},
 \quad c_4=-\frac5{128},\quad c_5=\frac7{256}.
\]
They are obtained by finite polynomial iteration, without selecting a
complex square root at each step. The root is the unique one in
\(\lVert Y-1\rVert_1\leq1/8\), and
\[
 \lVert P(x,1+S_K(x))\rVert_1\leq\tfrac14\,2^{-K}.
\]
The radius is intentionally conservative; the theorem prioritizes a
uniform finite proof over a sharp analytic radius.

For \(P(X,Y)=Y+Y^3-X\), take \(y_0=0\), \(d=1\),
\(Q(X,Z)=X-Z^3\), \(D=3\), and \(B=2\). The construction gives
\(s=1/24\), \(r=1/192\), and coefficients beginning
\[
  c_1=1,\quad c_2=0,\quad c_3=-1,\quad c_4=0,
  \quad c_5=3,\quad c_6=0,\quad c_7=-12.
\]
This is a concrete local inversion theorem with a fully specified
primitive recursive precision schedule.

Return now to the obstruction polynomial \(Y^2-X^2-a\). If a root
\(y_0\) of \(Y^2-a\) is supplied together with a separation witness
for \(2y_0\), the theorem gives its simple branch and exact
coefficient names. When \(a=0\), that witness does not exist for
\(y_0=0\), although the two branches \(Y=\pm X\) can be written
down in that special case. There is no uniform rule here for deciding
which case occurs. The theorem therefore resolves the regular branch
construction for Cauchy-presented coefficients while preserving the
repository's genuine obstruction at the singular transition.

\subsection{Finite mesh certificates replace one common use of compactness}
\label{sec:mesh-certificates}

The preceding result uses a supplied quantitative hypothesis. A related
repair applies whenever an argument invokes a positive minimum on a
compact domain. A finite mesh and a continuity modulus can supply the
needed lower bound, provided the finite numerical tests have enough
margin. They do not manufacture that margin from the bare assertion
that a function has no zeros.

\begin{proposition}[A finite certificate of positivity]
\label{prop:mesh-positivity}
Let \(f\) be a fixed real evaluator on a rational box
\([0,1]^d\), with an approximation procedure \(A(q,p)\) satisfying
\(|A(q,p)-f(q)|\leq e_p\) in the represented sense. Suppose a
primitive recursive modulus \(\mu\) certifies
\[
 \lVert x-y\rVert_\infty\leq2^{-\mu(k)}
       \quad\Longrightarrow\quad
 |f(x)-f(y)|\leq2^{-k}.
\]
Fix a positive rational \(\gamma\). Choose \(k,p\) so that
\(2^{-k}\leq\gamma\) and \(e_p\leq\gamma\), and let
\(M=2^{\mu(k)}\). If the finite tests
\[
 A(j_1/M,\ldots,j_d/M,p)\geq3\gamma
       \quad(0\leq j_i\leq M)
 \tag{A20}\label{eq:mesh-certificate}
\]
all hold, then \(f(x)\geq\gamma\) throughout the box.
\end{proposition}
\begin{proof}
For a rational point of the box, choose in each coordinate a nearest
grid point, or simply the lower grid point, with the endpoint \(1\)
handled by the last grid index. The sup-norm distance is at most
\(1/M\). Each tested grid value of \(f\) is at least
\(3\gamma-e_p\geq2\gamma\). The modulus then gives
\(f(x)\geq2\gamma-2^{-k}\geq\gamma\).
The grid indices and the selected point are obtained by bounded
integer division. The test (A20) is a bounded conjunction, and the
argument is a rational inequality after denominator clearing.
For a represented point, first clamp each coordinate of its rational
approximants to \([0,1]\); this operation is rational and does not
increase their distance from any point of the box. One can then apply
this rational theorem to the resulting approximants and the supplied modulus,
or search the finite grid against one approximant with an additional
precision margin. There is no need to decide which exact grid cell
contains an arbitrary real point.
\end{proof}

For a complex function, apply the proposition to a certified real
quantity such as \(|g|_2^2\). Polynomial data supply the continuity
modulus by finite difference estimates. A successful mesh certificate
then supplies the denominator bound required for inversion, root
separation, or a contour argument. An unsuccessful certificate proves
nothing about the presence of a zero. In particular, increasing the
mesh until a test succeeds is not justified as a primitive recursive
algorithm without a separately supplied success bound.

The same method yields a rational upper certificate for a supremum:
the largest grid approximation plus its evaluation error plus the
modulus error bounds every function value. A grid point attaining the
finite maximum is an approximate maximizer. This supplies the
quantitative consequence often used from an extreme-value theorem,
while keeping the required continuity and precision data explicit.

\section{Finite stability when no convergence modulus is available}
\label{sec:stability}

An explicit modulus of convergence is excellent input when it is available.
It should not, however, be made a prerequisite for every use of a bounded
sequence. Many arguments inspect only a finite interval of indices. For
such arguments a weaker, precisely quantified conclusion suffices. This
section gives one arithmetic selection principle that handles several
analytic observations and several finite exceptional sets at the same
time. The principle is a finite version of standard proof-mining ideas
\cite{ProofMining}; the additional purpose here is to expose exactly the
finite hypotheses and the quantifier-free invariants needed for iPRA.

\subsection{The finite data and the common-window theorem}

Fix a primitive recursive function $g:\NN\to\NN$ and an initial index $b$.
No monotonicity assumption on $g$ is required. Define
\begin{equation}\label{eq:endpoints}
 n_0=b,\qquad n_{i+1}=n_i+g(n_i)+1,
 \qquad W_i=[n_i,n_i+g(n_i)]\cap\NN.
\end{equation}
The windows are disjoint and ordered. All the endpoints are primitive
recursive, for a fixed primitive recursive definition of $g$.

For $j<m$, let $a_j(t)=a(j,t)$ for one fixed primitive recursive rational
evaluator $a$, let
$\epsilon_j>0$ be rational, and let $V_j\geq0$ be rational. For $\ell<r$,
let $E_\ell(t)=E(\ell,t)$ for one fixed primitive recursive Boolean
evaluator $E$, with natural
budget $C_\ell$. A successful window is one on which
\begin{align}
 |a_j(p)-a_j(q)|&\leq\epsilon_j
   &&(j<m,\ p,q\in W_i),\label{eq:stable-window}\\
 E_\ell(t)&=0 &&(\ell<r,\ t\in W_i).\label{eq:empty-window}
\end{align}
Additional natural parameters of these fixed evaluators are permitted;
the indices $j,\ell$ do not range over arbitrary program codes.
Every test in these formulas is finite and decidable. Use rational
cross-multiplication for the comparisons, not floating-point arithmetic.
Put
\begin{equation}\label{eq:stability-budget}
 K_j=\left\lfloor\frac{V_j}{\epsilon_j}\right\rfloor,
 \qquad K=\sum_{j<m}K_j+\sum_{\ell<r} C_\ell,
 \qquad T=n_{K+1}.
\end{equation}

\begin{theorem}[A common window from additive finite budgets]
\label{thm:common-window}
Suppose the following \emph{finite} inequalities hold:
\begin{align}
 \sum_{t<T}|a_j(t+1)-a_j(t)|&\leq V_j &&(j<m),\label{eq:finite-variation}\\
 \sum_{t<T} E_\ell(t)&\leq C_\ell &&(\ell<r).
 \label{eq:finite-exceptions}
\end{align}
Then a bounded search among $i=0,\ldots,K$ finds a successful window.
In particular its starting point is at most $n_K$. For fixed primitive
recursive evaluators $a_j,E_\ell,g$, this is an iPRA theorem with
quantifier-free finite input and output predicates.
\end{theorem}

\begin{proof}
For a fixed component $j$, call a window $j$-bad if the first test fails
for that component. Bounded search then supplies $p,q\in W_i$ with
$|a_j(p)-a_j(q)|>\epsilon_j$. Swapping $p,q$ if necessary and telescoping
gives
\[
 \epsilon_j<|a_j(p)-a_j(q)|
 \leq\sum_{t=p}^{q-1}|a_j(t+1)-a_j(t)|
 \leq\sum_{t=n_i}^{n_i+g(n_i)-1}|a_j(t+1)-a_j(t)|.
\]
For different $i$ these ranges of increments are disjoint. Consequently,
if $L$ windows are $j$-bad, then
$L\epsilon_j<\sum_{t<T}|a_j(t+1)-a_j(t)|\leq V_j$ when $L>0$.
There can therefore be at most $K_j$ such windows. This bound is slightly
generous when $V_j/\epsilon_j$ is an integer; generosity avoids a special
case and does no harm.

For a fixed exceptional predicate $E_\ell$, every window failing its
emptiness test contains a least exceptional index. The chosen indices
are distinct because the windows are disjoint, so there are at most
$C_\ell$ such windows by \eqref{eq:finite-exceptions}.

If all $K+1$ windows failed, assign each window to its least failed test,
in a fixed order of the $m+r$ tests. Each assigned class has size at most
its budget by the preceding two paragraphs. Summing the class counts
would give $K+1\leq K$, a contradiction. Because the tests are decidable,
this proves a bounded positive existence statement: scan the $K+1$
windows and retain the first success. If the scan returned failure,
the finite argument just given would contradict its recorded failures.
No double-negation elimination for an unbounded proposition is used.

All proof loops have visible finite bounds. The triangle inequality is
proved by induction on the number of increments; disjoint-range sums are
proved by induction on the number of selected windows; the last counting
identity is proved by induction on the list of assigned tests. Each
invariant is a rational inequality, a bounded conjunction, or an equality
of finite sums. After rational coding and Booleanization these are
quantifier-free formulas. The endpoint $T$ is fixed before these
inductions. Thus the proof has the required iPRA form.
\end{proof}

The use of variation through $T$ intentionally overestimates what is
needed: the last tested window ends at $T-1$. This harmless extra
increment makes the interface easy to compose. For $m=r=0$, $K=0$ and
the first window succeeds vacuously. A zero-length window also behaves
correctly: its oscillation is zero, while its exceptional-index test
still inspects its single index.

\subsection{Monotone sequences, finite variation, and simultaneous use}

If $a(t)$ is nondecreasing and $0\leq a(t)\leq M$, then for every finite
$T$ its variation telescopes to $a(T)-a(0)\leq M$. With
$\epsilon=2^{-k}$, Theorem~\ref{thm:common-window} therefore gives
\[
 \exists i\leq \lfloor M2^k\rfloor\quad
 \forall p,q\in[n_i,n_i+g(n_i)]\quad |a(p)-a(q)|\leq2^{-k}.
\]
The displayed quantifiers over $i,p,q$ are bounded; the search returns
the particular $i$. This is often called \emph{metastability}. Its
meaning is finite-window stability, not eventual stability on an
infinite tail.

The same theorem covers nonmonotone sequences with a proved finite
variation bound. It also synchronizes several sequences without first
finding a stable interval for one sequence and then searching inside it
for another. Their budgets add. Exceptional sets can be included in
the same selection. This matters when a proof needs, at one common
scale, a small oscillation in several approximants and avoidance of a
finite family of bad indices.

For example, suppose two monotone rational sequences take values in
$[0,1]$, both tolerances are $1/4$, and a decidable exceptional predicate
has at most two values in the inspected prefix. Then $K=4+4+2=10$.
Eleven candidate windows suffice for a common successful one. If
$g(n)=n+1$ and $b=0$, then $n_i=2^{i+1}-2$, so the returned start is
at most $2046$. This is a guaranteed bound; a particular input usually
succeeds much earlier.

\paragraph{Relation to the existing Roth packet.}
The repository already observes that a bound on the number of good
approximants yields a finite window with no such approximants
\cite{ProveIt}. Theorem~\ref{thm:common-window} does not claim that
observation as new, nor does it prove the missing quantitative Roth
estimate. It generalizes the common last step to multiple exceptional
predicates and analytic variation budgets, and supplies one precise
finite interface usable by other proofs. An absent cardinality estimate
remains an absent hypothesis.

\subsection{Why the output cannot be silently upgraded to a modulus}

The difference between a stable window and a Cauchy modulus is essential.
Let $H(e,t)$ be the decidable predicate that a fixed universal machine
with input $e$ has halted within $t$ steps. Bounded simulation is primitive
recursive. The sequence $a_e(t)=H(e,t)$ is nondecreasing and lies in
$[0,1]$. If a computable procedure uniformly returned a Cauchy cutoff
$M(e)$ for tolerance $1/2$, simulating through $M(e)$ would decide
halting: if the machine had not halted then, it could never halt later,
since that would change the sequence from $0$ to $1$ after the cutoff.
Thus even this very simple uniform family has no uniform computable
Cauchy modulus, although the finite-window theorem applies to it.

One can also see the obstruction in a single fixed primitive recursive
rational sequence. Define
\[
 s_t=\sum_{e\leq t}4^{-(e+1)}H(e,t).
\]
It is nondecreasing, primitive recursive, and bounded by $1/3$. Suppose
it had a computable Cauchy modulus. In the ordinary external real-number
interpretation, that modulus would compute its limit $s$. To recover the
halting status of machine $e$, first subtract the contributions of
machines $0,\ldots,e-1$, which have already been recovered. Put
$w=4^{-(e+1)}$. If $e$ never halts, the remaining value is at most $w/3$;
if it halts, the remaining value is at least $w$. An approximation with
error at most $w/6$ separates these cases, for example at threshold
$2w/3$. Recursing through $e$ decides the halting problem, a contradiction.
This establishes that a global modulus cannot in general be extracted
even for a single primitive recursive bounded monotone rational sequence.

The external impossibility argument is not used as an axiom in any
positive iPRA proof. It diagnoses a proposed change of output contract.
Nor does it license a blanket claim that an arbitrary mathematical
theorem is unprovable in iPRA. One must first identify its exact input
representation and exact conclusion.

\subsection{The Lean and arithmetic boundaries of the selector}

The shipped \texttt{FiniteCore.lean} defines the bounded scan
\texttt{firstHit}, disjoint endpoints, and finite contracts. It proves
that a successful index within the budget makes the scan succeed, that
the scan stays within the budget, and that bounded contracts compose.
It also proves a natural-valued energy termination lemma, a useful
alternative implementation when the failed-test budget is already an
integer potential.

These are compiled ambient Lean theorems. The rational variation and
multi-component counting proof above has a formalization outline, not a
completed Lean proof. To finish it, represent the rational arrays by
finite lists; prove interval telescoping, the disjoint-sum inequality,
and the finite assignment count; and then apply
\texttt{firstHit\_complete}. To finish the \emph{object-theory} proof,
replace the arrays by natural codes and construct the corresponding
\texttt{Prf PRAAx} derivations using the bounded invariants just listed.
Compilation of the ambient selector alone does not perform that last
step.

\section{Repository difficulties and a concrete repair in Diophantine geometry}
\label{sec:repository-repair}

\subsection{What the inspected packets actually leave open}

The audit in this section refers throughout to ProveIt commit
\path{27a07921176bccbf630084b959700605783ff4aa}.
The index under \path{Logic/PrimitiveRecursiveArithmetic/docs/proofs/}
lists 47 delivered packets and explicitly says that their placement did
not rerun, compile, verify, or edit them.  A packet's own claim of a
compiled native Lean theorem must therefore be distinguished both from an
independent build in this project and from an object-language iPRA
derivation.  We use the packets as inspectable mathematical arguments and
interfaces, rather than treating their index as a correctness certificate
\cite{ProveIt}.

Several obstacles in the packets are already identified with considerable
precision.  It would not advance the program to rename their existing
solutions.  In particular, the Mordell--Weil packet already has contracting
height descent, logarithmic fuel, explicit coordinates, and a method of
obtaining a bounded selector from actual covering representatives.  The
Siegel packet already has a finite-cover-to-cardinality lemma and observes
that covering subspaces of a supplied tuple can be replaced by spans of
subtuples.  Neither result is claimed here as a new repair.

\begin{center}
\small
\begin{tabular}{p{0.19\linewidth}p{0.36\linewidth}p{0.36\linewidth}}
\hline
Packet & Identified missing information & Consequence for this framework\\
\hline
\texttt{puiseux\_ipra} & Arbitrary fast Cauchy names do not decide the
exact coefficient branches needed for full Puiseux output. & Preserve the
exact algebraic input interface; approximation alone does not supply
its zero tests.\\[3pt]
\texttt{mordell\_weil\_ipra} & A finite list of actual rational
representatives for the quotient is still required. & The existing descent
algorithm can consume the list, but does not manufacture it from Selmer
information.\\[3pt]
\texttt{faltings\_ipra} and \texttt{siegel\_ipra} & General positive
height or complete-list bounds are not supplied. & A count bound cannot
serve as a search cutoff.  We prove a useful count ingredient below,
without claiming a height bound.\\[3pt]
\texttt{diophantine\_ipra} & The general effective Roth cutoff is not
proved; quadratic and Liouville-range cases are treated. & Keep the
denominator cutoff distinct from the number of exceptional approximants.\\[3pt]
\path{prime_distribution_ipra} & Chebotarev needs certified
field arithmetic and quantitative analytic estimates; full SW and BV
need the indicated uniform witnesses. & Finite prime enumeration and
finite-precision logarithms do not provide those estimates.\\
\hline
\end{tabular}
\end{center}

These descriptions are claims about the inspected reports, not an
assertion that every stronger theorem is unprovable in iPRA.  Their exact
source paths and blob identifiers are recorded in the accompanying source
manifest.  The following repair targets a specific obligation in
\path{siegel_ipra/paper/05_general_route.tex}: the bounded number of curve
points lying over a proper subspace.  That report sketches a divisor
argument but leaves its arithmetic realization open.  For odd-degree
hyperelliptic equations, the realization can be given directly by one
polynomial norm, with an explicit finite proof.

\subsection{A zero-count theorem from a polynomial norm}

All polynomials and coordinates in this subsection are rational and
finitely encoded.  Rational equality and order mean exact comparisons
after cross multiplication; they are not comparisons of Cauchy names.
Let $f\in\mathbb Q[X]$ have positive odd degree $d$, with nonzero leading
coefficient.  Let
\[
             C_f=\{(x,y)\in\mathbb Q^2:y^2=f(x)\}.
\]
No smoothness assumption is needed for the arithmetic theorem.

For $A,B\in\mathbb Q[X]$, not both zero, define
\[
 G(x,y)=A(x)+B(x)y,\qquad
 N(X)=A(X)^2-f(X)B(X)^2.
\]
In degree inequalities, a zero polynomial simply contributes no term:
thus $2\deg A\le L$ is required only when $A\ne0$, and
$d+2\deg B\le L$ only when $B\ne0$.

\begin{theorem}[Arithmetic fiber bound]
\label{thm:hyperelliptic-fiber}
The polynomial $N$ is nonzero.  Every finite list of pairwise distinct
points of $C_f$ satisfying $G(x,y)=0$ has length at most $\deg N$.
Consequently it has length at most $L$ whenever
\[
        2\deg A\le L,\qquad d+2\deg B\le L
\]
for the nonzero terms.  The theorem has a proof using rational polynomial
arithmetic, bounded finite-list operations, and quantifier-free induction.
\end{theorem}

\begin{proof}
First suppose $B=0$.  Then $A\ne0$ and $N=A^2\ne0$.  If $A=0$ and
$B\ne0$, then $N=-fB^2\ne0$.  If both are nonzero, the degrees of $A^2$
and $fB^2$ have different parity: the former is even and the latter odd.
Their highest terms therefore cannot cancel.  This proves nonvanishing,
and also
\[
 \deg N\le \max\{2\deg A,d+2\deg B\},
\]
with absent terms omitted.

Now take a finite pairwise-distinct list of common zeros.  Group it by
its $x$-coordinate using exact rational equality.  For each different
$x$-coordinate $a$, write $m_a$ for the number of points in that group.
The equation $y^2=f(a)$ has at most two distinct rational solutions:
if $u^2=v^2$, then $(u-v)(u+v)=0$, and the rational field has no zero
divisors.  Thus $m_a\le2$.

Every nonempty group gives $N(a)=0$, by substituting
$A(a)=-B(a)y$ and $y^2=f(a)$.  Hence $X-a$ divides $N$.
If $m_a=2$, let the two different ordinates be $u,v$.  Subtracting the
two equations $A(a)+B(a)u=0$ and $A(a)+B(a)v=0$ gives
$B(a)(u-v)=0$.  Since $u-v\ne0$, we have $B(a)=0$, and then $A(a)=0$.
The factor theorem yields $A=(X-a)A_1$ and $B=(X-a)B_1$, so
\[
          N=(X-a)^2(A_1^2-fB_1^2).
\]
This argument includes $B=0$ and handles branch points as well: a branch
point has only one ordinate and requires only one factor.

It follows that
\[
              \prod_a(X-a)^{m_a}\ \text{divides}\ N.
 \label{eq:weighted-roots}
\]
For completeness, this simultaneous divisibility uses no factorization
oracle.  Process the distinct $a$'s in their list order.  Suppose the
product $C$ of the already processed factors has been removed, so
$N=CQ$.  At the next coordinate $a$, every factor of $C$ has nonzero
value at $a$; hence $C(a)\ne0$.  The equality $N(a)=0$ gives $Q(a)=0$,
and synthetic division removes one $X-a$ from $Q$.  When $m_a=2$,
the known divisibility of $N$ by $(X-a)^2$, followed by cancellation of
one $X-a$, gives a second vanishing at $a$ and permits a second division.
All cancellations are by nonzero polynomials in $\mathbb Q[X]$.
This proves the displayed product divisibility by a finite induction.

The product is monic and has degree $\sum_a m_a$, which is exactly the
length of the original list.  Since $N\ne0$, division cannot increase
degree.  Therefore
\[
           \#\text{points}=\sum_a m_a\le\deg N\le L.
\]
Every operation just described has an explicit finite implementation.
Polynomial division by $X-a$ is synthetic division along the coefficient
list, and testing $A=0$ or $B=0$ is a finite coefficient check.
\end{proof}

\paragraph{Why counting only distinct $x$-coordinates loses information.}
A rough argument would count roots of $N$ and then multiply by two.  The
theorem avoids that loss because a fiber containing two common zeros
already consumes two factors of $N$.  This improvement is attained.
Take
\[
       f=X^3-X+1,\qquad A=X(X-1),\qquad B=0.
\]
Then $N=X^2(X-1)^2$ has degree four.  Its common zeros on the curve are
exactly
\[
                  (0,1),\ (0,-1),\ (1,1),\ (1,-1).
\]
Indeed $A(x)=0$ forces $x=0$ or $x=1$, and in either case $f(x)=1$.
Thus the bound four is exact; the rough $2\deg N$ bound would be eight.

\subsection{An explicit replacement for the divisor-based fiber contract}

Fix $L\ge d$ and put
\[
 a=\lfloor L/2\rfloor,\qquad b=\lfloor(L-d)/2\rfloor,\qquad
 n=a+b+2.
\]
Use the polynomial evaluation vector
\[
 \Phi_L(x,y)=(1,x,\ldots,x^a,\ y,xy,\ldots,x^b y)
                      \in\mathbb Q^n.
\]
The weight two assigned to $x$ and the weight $d$ assigned to $y$
are only a convenient organization of the norm degree in the proof;
no theory of divisors is being assumed.

\begin{corollary}[Proper subspaces have uniformly bounded fibers]
\label{cor:subspace-fiber}
For every proper rational linear subspace $W\subsetneq\mathbb Q^n$,
at most $L$ distinct rational points of $C_f$ have $\Phi_L(x,y)\in W$.
More precisely, every such finite list has length at most $L$.
\end{corollary}

\begin{proof}
Represent $W$ by a finite matrix of spanning vectors and the assertion
that its row rank is less than $n$.  Rational Gaussian elimination finds
a nonzero vector in its annihilator: select a free column, assign its
unknown the value one, assign the other free unknowns zero, and solve
backwards for the pivot unknowns.  This is a bounded exact procedure.
Write the resulting linear form on $\Phi_L$ as $A(x)+B(x)y$.
Its coefficient vector is nonzero, so $A,B$ are not both zero, and
its displayed degree bounds imply the hypotheses of
Theorem~\ref{thm:hyperelliptic-fiber}.  Every point mapping into $W$
is a zero of this form.  Apply that theorem.
\end{proof}

\begin{corollary}[A finite-tuple approximation theorem suffices]
\label{cor:hyperelliptic-cardinality}
Suppose a proved arithmetic contract says that for every finite list
of distinct points in a specified decidable subset $S\subseteq C_f$,
the corresponding vectors $\Phi_L$ can be covered by at most $M$ proper
rational subspaces.  Then every finite list of distinct points in $S$
has length at most $ML$.
\end{corollary}

\begin{proof}
Given the finite covering list, assign each point to the first subspace
containing its vector.  Membership is a finite rational rank test, and
coverage makes the bounded scan succeed.  Each of at most $M$ groups
contains at most $L$ points by Corollary~\ref{cor:subspace-fiber}.
Adding the group sizes gives the bound $ML$.  Equivalently, applying
this argument to a hypothetical list of $ML+1$ distinct points yields
an arithmetic contradiction.
\end{proof}

This explicitly discharges the zero-count half of the proposed Siegel
strategy for the odd-degree hyperelliptic evaluation systems above.
It does not prove the quantitative approximation theorem needed for the
covering contract, and hence does not prove general Siegel or Faltings.
The result supplies a cardinality bound once that contract is established;
it does not locate the points or give a coordinate cutoff.  The repair is
nevertheless substantive: one entire geometric dependency, previously
expressed through Riemann--Roch functions and zero/pole divisors, is
replaced here by a proved finite arithmetic lemma on this class.

\subsection{Finite algebra certificates as a general zero-count interface}

The preceding norm argument is particularly efficient.  A more general
interface works for arbitrary finite polynomial systems and can be reused
where a small norm formula is unavailable.

Let $F_1,\ldots,F_s\in\mathbb Q[X_1,\ldots,X_r]$.  A
\emph{spanning certificate of size $D$}, with $D\ge1$, consists of
polynomials $b_1=1,b_2,\ldots,b_D$, rational numbers $m_{kji}$, and
polynomials $u_{kj\ell}$ satisfying the finitely many identities
\[
 X_k b_j-\sum_{i=1}^D m_{kji}b_i
       =\sum_{\ell=1}^s u_{kj\ell}F_\ell
 \quad(1\le k\le r,\ 1\le j\le D).
 \label{eq:spanning-certificate}
\]
The checker expands and compares polynomial coefficients.  It checks
fixed algebraic identities, not the truth of arbitrary formal proofs.
Linear independence of the $b_i$ is not required; the certificate gives
an upper bound on dimension and can therefore be deliberately redundant.

\begin{theorem}[Certificate bound on rational zeros]
\label{thm:finite-algebra-zeros}
If a size-$D$ spanning certificate is valid, every finite list of distinct
common rational zeros of the $F_\ell$ has length at most $D$.
The checker, the reduction operations, and the finite contradiction
procedure used in the proof are primitive recursive.
\end{theorem}

\begin{proof}
The identities imply, by multiplying one variable at a time, that every
polynomial is congruent modulo $(F_1,\ldots,F_s)$ to a rational linear
combination of the $b_i$.  One can implement this reduction without a
Gr\"obner-basis algorithm: start with the coefficient vector of $1$,
apply the relevant matrix $m_k$ for each variable occurrence in a
monomial, and add the resulting vectors with the monomial coefficients.
For correctness, substitute a common zero $p$ in each certificate
identity.  Induction on the finite multiplication sequence shows that
the resulting coefficient vector $c$ satisfies
\[
                   H(p)=\sum_i c_i b_i(p)
\]
for the polynomial $H$ being reduced.  This equality is all that the
argument needs.

Suppose distinct common zeros $p_1,\ldots,p_t$ have been supplied.
For each ordered pair $i\ne j$, scan the $r$ coordinates to choose
the first $k(i,j)$ where $p_i$ and $p_j$ differ.  Their distinctness
makes the scan succeed.  Form
\[
 H_i(X)=\prod_{j\ne i}
 \frac{X_{k(i,j)}-(p_j)_{k(i,j)}}
      {(p_i)_{k(i,j)}-(p_j)_{k(i,j)}}.
\]
Every displayed denominator is a witnessed nonzero rational.  Direct
substitution gives $H_i(p_j)=1$ if $i=j$, and zero otherwise.
Reduce $H_i$ to a coefficient vector $c_i\in\mathbb Q^D$ by the
certificate matrices.  If a linear relation
$\sum_i\lambda_i c_i=0$ holds, evaluation at $p_j$ gives
\[
       0=\sum_i\lambda_i H_i(p_j)=\lambda_j.
\]
Thus the $t$ coefficient vectors are linearly independent.

Here the last dimension step is an explicit finite one.  If $t>D$,
perform rational row reduction on the $t$ by $D$ matrix with rows $c_i$,
carrying along the same row operations on the $t$ by $t$ identity matrix.
There are at most $D$ pivot rows, so at least one resulting coefficient
row is zero.  The corresponding row of the companion matrix is nonzero,
because every elementary row operation is invertible.  It gives a
nonzero rational relation among the $c_i$, contrary to the preceding
evaluation argument.  This proves $t\le D$.

All loops are bounded by the input lengths, numbers of variables,
exponents, and matrix dimensions.  Arithmetic may create large rational
coefficients, but addition, multiplication, normalization, and division
by the selected nonzero pivots are primitive recursive on rational codes.
No unbounded search for a basis, an algebraic closure, or a solution occurs.
\end{proof}

The finite common-zero theorem works uniformly in the encoded polynomial
system and in the certificate.  This is legitimate uniformity: polynomial
evaluation and coefficient comparison are specific primitive recursive
algorithms.  It is not a universal evaluator for arbitrary primitive
recursive programs.

\paragraph{A size-four certificate for the sharp example.}
For
\[
 F_1=y^2-x^3+x-1,\qquad F_2=x^2-x,
\]
take $b=(1,x,y,xy)$.  The elementary identities
\[
        x^2=x+F_2,\qquad
        y^2=1+F_1+(x+1)F_2
\]
and their multiples by $x$ or $y$ give all eight required border
identities.  Multiplication by $x$ sends the four basis entries, modulo
the two equations, to $(x,x,xy,xy)$; multiplication by $y$ sends them
to $(y,xy,1,x)$.  Therefore the certificate theorem gives at most four
rational solutions, and the four displayed points show that this bound
is attained.  The coefficient identities themselves can be checked by
an exact polynomial program without any analytic or geometric premise.

One should separate certificate \emph{verification} from certificate
\emph{construction}.  A system can have infinitely many zeros and then
admit no such certificate.  Even for a system known to be zero-dimensional,
calling an unspecified Gr\"obner-basis implementation does not by itself
prove a primitive recursive bound or an iPRA correctness theorem for its
construction.  Supplying a concrete certificate, or proving a bounded
constructor for the intended family, is the remaining input.

\subsection{An induction audit and Lean completion path}

Both zero-count arguments have particularly simple object-theory
invariants.  For Theorem~\ref{thm:hyperelliptic-fiber}, encode a finite
point list, its distinct-coordinate groups, a residual polynomial, and
the product of the factors already removed.  The loop invariant is the
conjunction of polynomial coefficient equalities
\[
 N=CQ,\qquad
 \deg N=\deg C+\deg Q,
\]
the accumulated root multiplicity, and the finite validity checks on the
processed groups.  It is used only while $N,Q$ are nonzero.  Nonzero
coefficient tests, polynomial degrees, list membership, and bounded
universal tests are all represented by primitive recursive Booleans.
Synthetic division carries its own coefficient identity as a second
finite invariant.  Consequently the induction formula contains no
unbounded quantifier over points or polynomials.  The arbitrary input
point list is a parameter.

For Theorem~\ref{thm:finite-algebra-zeros}, the reduction invariant is
evaluation equality for all entries of the supplied finite point list,
expressed by one bounded Boolean fold.  Row reduction carries the
coefficient matrix, the companion matrix, and its inverse; the equations
that these two square matrices multiply to the identity prove that the
companion row used at the end is nonzero.  Pivot selection is a bounded
scan with an explicit failure branch.  A failed pivot means the column
has no candidate, rather than an appeal to choice.  These are again
quantifier-free equalities on coded finite objects.

The formalization can be organized in the following order.
\begin{enumerate}
\item Define canonical rational codes, polynomial coefficient lists, finite
multivariate monomial lists, their primitive recursive evaluation
programs, and Boolean equality tests.  Prove the field and polynomial
identities used by synthetic division and matrix arithmetic in the
object theory.
\item Prove the synthetic-division identity and its specialization to a
root.  Prove nonvanishing and degree of products from leading coefficients.
Implement the distinct-coordinate grouping procedure and its
length-preservation identity.
\item Formalize the multiplicity-two fiber lemma: two distinct ordinates
annihilated by $A+B y$ force $A=B=0$ at that abscissa.  Combine its
explicit factors with the root-removal invariant to obtain the
hyperelliptic bound.
\item Define the spanning-certificate Boolean by coefficientwise checks.
Prove the reduction invariant, construct the finite separator vectors,
and prove their evaluation matrix is the identity.
\item Formalize bounded Gaussian elimination, including the tracked
inverse for row operations.  Use it for both a proper-subspace
annihilator and the finite-algebra dimension contradiction.
\item Assemble actual embedded derivations of the zero-count formulas.
Only after these are constructed should the existing ConciseIPRA
interpretation machinery transport them to another presentation of the
theory.
\end{enumerate}

For example, the native specification of the final fiber statement has
the following shape, where every predicate in the hypotheses is a
decidable finite-code predicate and \texttt{points.length} refers to an
encoded list length:
\begin{verbatim}
def HyperellipticFiberBound : Prop :=
  forall f A B points,
    OddDegree f -> (Nonzero A || Nonzero B) = true ->
    PairwiseDistinct points = true ->
    AllCurveAndLinearZeros f A B points = true ->
    points.length <= degree (A*A - f*B*B)

def FiniteAlgebraZeroBound : Prop :=
  forall equations certificate points,
    CheckSpanning equations certificate = true ->
    PairwiseDistinct points = true ->
    AllCommonZeros equations points = true ->
    points.length <= certificate.dimension
\end{verbatim}
This pseudocode specifies the mathematical types; it is not presented as
a compiled Lean module.  In an actual implementation, the polynomial
operators, code projections, and Boolean names must be replaced by the
chosen concrete definitions.  The repository endpoint is a corresponding
object-language formula together with a term of type
\texttt{Prf PRAAx [] formula}, not merely a proof of the native
\texttt{Prop}.  Neither a field in a Lean structure nor an uninstantiated
contract should be counted as that derivation.


\section{Lean formalization: definitions, targets, and the completion path}
\label{sec:lean-plan}

The formalization has three distinct endpoints. First, executable native
definitions fix the finite objects and prevent an informal theorem from
changing shape during implementation. Second, native correctness lemmas
establish the finite mathematics in Lean. Third, embedded derivations
establish the same assertions in the specified iPRA calculus. This report
provides a small compiled part of the second endpoint, compiled definitions
for several of the remaining targets, and a detailed plan for the third.

\subsection{The existing repository interfaces}

The pinned repository's
\path{ConciseIPRA/Lean/ConciseIPRA/PR.lean} defines an inductive type
\texttt{PR k} with constructors for zero, successor, projection,
composition, and primitive recursion. It also defines \texttt{PTm n},
the terms over $n$ free variables. The syntax file supplies
\texttt{Fm PTm n}; the deduction file supplies the derivation type
\texttt{Prf Ax Gamma phi} for intuitionistic natural deduction with
equality. The arithmetic axiom family is \texttt{PRAAx}.

In particular, the repository's induction constructor accepts a proof
that its formula is quantifier-free. The derived rule
\texttt{PRA.qfInd} takes the base derivation and a derivation of the
successor step under the shifted context. A completed theorem in the
present report should end with a declaration of the following shape:
\begin{verbatim}
def target : Fm PTm parameterCount := ...
theorem target_derivation : Prf PRAAx [] target := ...
\end{verbatim}
The ellipses here denote implementation work explicitly listed below;
they are not part of a claimed compiled theorem.

The module \texttt{NativeBiInterpretation.lean} supplies
\texttt{pra\_iff\_native}. It transports an already supplied iPRA
derivation to the native presentation, using \texttt{NFm.ofPRA}.
It does not turn an arbitrary Lean proposition into an iPRA theorem.
The archive's \texttt{ObjectInterface.lean} records these actual types,
fixed PR name symbols, syntactic implication composition, and the
transport step. That integration file has not been compiled here,
because it requires the full repository library and Mathlib. Its status
is separate from the three standalone compiled modules.

\subsection{The compiled native definitions}

\paragraph{Finite search and contracts.}
\texttt{FiniteCore.lean} uses only Lean's standard distribution. It defines
a bounded scan \texttt{firstHit p n} that retains the first successful
index in $0,\ldots,n$. If no success exists it returns $n$; therefore
success is always tested by applying \texttt{p} to the returned value,
not inferred from the index alone. This convention differs from the
sentinel implementation in Lemma~\ref{lem:bounded-choice}, but both
are explicitly bounded and their success conditions are unambiguous.

The file proves the bound on the returned index, completeness from a
supplied successful index, and the failure certificate for every
inspected index. Its \texttt{energy\_firstHit} lemma proves success
whenever a natural-valued energy starts below $B$ and drops by at least
one on each failed test through $B$. The proof maintains
$E(i)+i\le B$ under a hypothetical all-failure prefix. At $i=B+1$
that is impossible. The file also defines disjoint endpoints and
proves their monotonicity, composes finite contracts, and packages the
selected-window bound.

\paragraph{Polynomial and observation targets.}
\texttt{Targets.lean} fixes concrete native data:
\begin{itemize}
\item rational polynomials as coefficient lists, with coefficient lookup,
evaluation, addition, multiplication, degree, nonzero tests, and weighted
finite coefficient sums;
\item rational point lists, exact pairwise distinctness, the common-zero
test for $y^2=f(x)$ and $A(x)+B(x)y=0$, and the norm polynomial;
\item finite lists of bivariate terms, coefficient collection, substitution,
and polynomial iteration truncated modulo $X^{K+1}$;
\item regular rational-name predicates, finite regularity tests, finite
variation, and finite window-stability tests;
\item bivariate spanning-certificate data, matrix rows, ideal-multiplier
witnesses, dimension checks, exact border-identity checks, and the
common-zero predicate.
\end{itemize}
Definitions named \path{HyperellipticFiberTarget},
\path{ImplicitCoefficientStabilityTarget},
\path{FiniteVariationTarget}, and \path{FiniteAlgebraTarget}
are propositions specifying the remaining theorems. They are \emph{not}
theorems or assumed axioms. The bivariate certificate target is a concrete
two-variable instance of the general $r$-variable theorem in the paper.
Its checker verifies that the basis is nonempty and begins with $1$,
that all matrix and witness rows have the correct lengths, and that
every required polynomial identity holds after coefficient collection.

The finite-variation target uses the sufficient hypothesis
\[
 \operatorname{variation}(T)<(B+1)\epsilon
\]
to conclude that one of $B+1$ windows succeeds. This matches the
floor-budget construction in Section~\ref{sec:stability}. The native
functions in these statements can be arbitrary host functions. Their
additional fixed-PR representation is required when constructing the
object-theory schema; the native type alone does not certify it.

The source convention for the degree of the zero polynomial is zero.
This is harmless in the fiber target because its hypotheses imply
that the norm polynomial is nonzero. Polynomial multiplication may
retain trailing zero coefficients; coefficient comparison and the
implemented degree ignore those zeros. Weighted-norm lemmas must assume
their weight is nonnegative, and the analytic lemmas in the paper use
a strictly positive rational weight.

\paragraph{The Gaussian-rational branch prototype.}
\texttt{FiniteBranch.lean} adds the paired-coordinate version of the
polynomial algorithms and records the finite analytic goals in a
\texttt{KernelGoals} structure. No instance of that structure is supplied.
It also checks five individual coefficients of the branch of $Y^2=1+X$
by kernel computation. The companion \texttt{ANALYTIC\_PLAN.md}
specifies the finite norm, perturbation, and name proofs still required.

\subsection{Constructing the embedded arithmetic layer}

The following sequence gives concrete milestones. Each milestone has
a finite correctness statement and an indicated proof method.

\paragraph{1. Natural coding and arithmetic.}
Choose one total PR pairing function with bounded projections and a
finite-list code with PR length and entry functions. A radix list code
with an explicitly stored base is sufficient whenever a finite maximum
of the entries is available; a standard nested pairing is sufficient
otherwise. Integers may be represented by pairs of naturals, with
equality $(a,b)\simeq(c,d)$ defined by $a+d=b+c$. Represent rationals by
an integer numerator and a positive denominator. Do not assume that
arbitrary host \texttt{Rat} computations are already object-theory
functions.

Build the PR definitions and derivations for arithmetic, comparison,
bounded sums, bounded maxima, and list operations. Prove validity
preservation separately. Cross-multiplication turns the rational
inequalities needed below into integer inequalities and then PR Boolean
tests. No quotient construction is needed inside iPRA.

\paragraph{2. Bounded formulas and proof construction.}
Implement the formula-to-Boolean translation of
Lemma~\ref{lem:bounded-booleanization} as a metalevel transformation
on a fixed formula syntax, emitting both a PR definition and its
equivalence derivation. Implement the bounded search derivations next.
The syntactic compiler then uses substitution, intuitionistic logical
rules, and \texttt{PRA.qfInd} on the emitted Boolean invariant. In
particular, it must never invoke ordinary Lean induction on an arbitrary
semantic proposition and declare the resulting proof to be an embedded
derivation.

\paragraph{3. Finite algebra.}
Prove coefficient-array convolution, distributivity, the difference of
powers identity, and synthetic division. The normal-form soundness
argument in Proposition~\ref{prop:polynomial-normalization} gives a
reusable source of finite identity derivations. Its ring operations are
fixed PR definitions with supplied arithmetic laws. Parameterized matrix
identities still need their uniform generators or direct proofs.

\paragraph{4. Analytic finite kernels and names.}
Use paired rational coordinates for Gaussian rationals. Prove their
$\ell^1$ triangle and product bounds. Prove (A1)--(A9) on finite
polynomial arrays; the induction invariants are the weighted norm bound,
the bounded list of zero low coefficients, and the rational perturbation
inequality. Fix the requested coefficient degree and both precision
indices before induction. Then instantiate these finite lemmas at the
explicit indices (A12)--(A17). The input regularity and root equation are
universal arithmetic assumptions that are used at the finite numerical
instances demanded by a particular conclusion.

For represented equality, adopt one fixed convention. For example, two
regular names $a,b$ represent the same value when for every requested
$n$ their sufficiently reindexed approximants differ by at most
$2^{-n}$. State the exact reindexing in each operation and verify its
transitivity by a further fixed precision shift. No equality decision
procedure is introduced. For root uniqueness, use the finite
contraction-with-error inequality and choose its error below the
requested tolerance.

\paragraph{5. Finite stability.}
Prove telescoping and the triangle inequality for a supplied finite
sequence, then the disjoint-increment sum bound. A failed window
contains a boundedly found witness pair. The finite assignment table
records which component pays for each failed window. Its invariant is
that all assigned indices lie in the processed prefix, that the classes
partition that prefix, and that each count respects its proved budget.
This produces the successful index needed by the already compiled
native search theorem, and separately produces an embedded proof of
the coded selector's correctness.

\paragraph{6. Hyperelliptic and finite algebra counts.}
For the norm theorem, carry the point groups, the removed product, and
the residual polynomial. The central invariant is $N=CQ$ together with
the degree sum and the number of points accounted for. The multiplicity
two lemma is finite rational field arithmetic. For the certificate
theorem, carry reduction vectors and the evaluation matrix of the
separator polynomials on the supplied point list. Bounded rational row
reduction tracks an invertible companion matrix, as described in
Section~\ref{sec:repository-repair}. These invariants are Boolean
equalities over finite codes, so their induction is permitted.

\paragraph{7. Assembly and audit.}
Construct the concluding \texttt{Prf PRAAx} terms. Audit every premise:
the polynomial identities, root equation, derivative gap, continuity
modulus, finite variation bound, spanning identities, or finite covering
theorem must either be input data with the stated contract or have an
actual preceding derivation. Only after this assembly invoke
\texttt{pra\_iff\_native}. There is no need to reprove bi-interpretability
for every application.

\subsection{What the formalization does and does not certify}

The lemmas in \texttt{FiniteCore.lean} prove native finite-search and contract
facts. Their printed axiom dependencies include Lean's standard
\texttt{propext} and \texttt{Quot.sound} for the arithmetic tactic
proofs; the contract-composition proof has no axiom dependencies.
They contain no \texttt{sorry}, no new axioms, and no use of
\texttt{Classical.choice}. This does not by itself certify that these
native proof terms belong to iPRA. The explicit arithmetic invariants
and the planned construction of \texttt{Prf PRAAx} terms are what address
that additional requirement.

The five computation proofs in \texttt{FiniteBranch.lean} have a
different printed dependency set: it includes \texttt{Classical.choice}
as well as \texttt{propext} and \texttt{Quot.sound} through their
native rational/container infrastructure. They are disclosed native
metatheory computations. The intended object proof instead uses the
raw integer/denominator representation and finite arithmetic arguments
given in the report; these computation proofs are not imported as
unexamined iPRA axioms.

The report's human proofs are substantially more developed than the
current Lean implementation. The remaining work is finite and specified
for the proved lemmas, but it includes real implementation and proof
construction: no blanket promise that a future \texttt{ring} or
\texttt{simp} call will discharge it is being made. The general
Diophantine approximation and height theorems identified in the audit
are further mathematical problems, not merely this implementation work.

\section{Validation, scope, and a practical continuation order}

\subsection{What was checked in this package}

The standalone Lean files \texttt{FiniteCore.lean},
\texttt{Targets.lean}, and \texttt{FiniteBranch.lean} were compiled
with Lean 4.32.0. The first
contains actual proofs of the bounded search, energy, endpoint, and
contract lemmas. The second fixes concrete definitions and propositions
for future proofs. Its successful compilation establishes that those
definitions and statements are well-typed, not that the target
propositions have been proved. The archive retains the compiler output
and its scope declaration. The third supplies Gaussian-rational
algorithms and five kernel-checked coefficient computations; its
general analytic goal fields remain uninstantiated.

The standard-library Python reference program uses exact
\texttt{Fraction} arithmetic and checks the following.
\begin{itemize}
\item All $256$ binary sequences of length eight, extended constantly,
are fed to the finite-variation selector with their actual variation
budget. Nonmonotone examples are included. Two inputs violating the
finite budget hypotheses are rejected. A two-sequence example with an
exceptional set verifies the synchronized output and the bound $2046$.
\item The simple branch of $Y^2=1+X$ is constructed through order $20$
by truncated polynomial iteration. Its coefficients agree with the
independent binomial recurrence. The run checks $21$ weighted bounds,
$231$ stabilization comparisons, $924$ finite-tail inequalities, and
$84$ residual inequalities, all exactly.
\item The sharp four-point hyperelliptic example and a second example
with $B\ne0$ verify the curve equations, norm polynomial, repeated
synthetic divisions, and point count.
\item All eight border identities for the size-four quotient certificate
are verified by coefficient comparison. The evaluation matrix on the
four displayed solutions has determinant $-4$, confirming that the
example attains the dimension bound.
\end{itemize}
These are implementation and example checks. The proofs of the universal
mathematical assertions are in the manuscript; their full machine
formalization remains as described above. In particular, no numerical
test has been used as evidence for an unresolved global height estimate.

The mathematical chapters were also reviewed independently during
preparation. The review checked the degree and multiplicity argument,
the finite algebra dimension contradiction, the common-window endpoint
conventions, the analytic norm and precision constants, and the
quantifier-free induction boundaries. The report was compiled to PDF
and its rendered pages inspected. The archive records reproducibility
commands, source provenance, and file checksums.

\subsection{Unresolved mathematical dependencies}

The proposed framework does not establish the following stronger claims:
\begin{enumerate}
\item a uniform primitive recursive bound for all rational points of
arbitrary genus-at-least-two curves, or for integral points on arbitrary
curves covered by Siegel's theorem;
\item a uniform generator algorithm for arbitrary Mordell--Weil groups
from abstract finiteness or Selmer information alone;
\item the finite-tuple quantitative Subspace Theorem required by the
conditional hyperelliptic point-count corollary;
\item a general effective Roth cutoff, or all of the analytic estimates
and uniform witnesses left open by the prime-distribution packet;
\item uniform exact singular Puiseux coefficient selection from arbitrary
Cauchy presentations, or a generic zero test for such presentations;
\item a global Cauchy modulus extracted from a finite-window selector;
\item an implemented automatic compiler from arbitrary Lean proofs to iPRA.
\end{enumerate}
The first four entries are left open here. Their absence is not a proof
of unprovability. The fifth and sixth contain precise representational
or computability obstructions, explained in the inspected source and
in this report. The last entry would require much more than the
restricted proof-construction discipline established here.

\subsection{Recommended implementation order}

The best first complete embedded theorem is the hyperelliptic fiber
bound. Its data are exact rationals, its new idea is isolated, and its
induction invariants are short. It also produces an immediate usable
contract for an existing research route. The size-four example supplies
a small end-to-end test of both the norm and the finite algebra paths.

Next implement polynomial certificate soundness and the spanning
certificate theorem. These modules can serve a broad range of algebraic
and geometric packets. In parallel, the bounded selector and
common-window theorem can be translated into the object theory using
only finite sums and counting.

The analytic branch theorem should then reuse the same finite polynomial
arithmetic, adding weighted norms and precision contracts. Its key
advantage is that the perturbation bound is independent of iteration
count. This provides a stable interface for later analytic applications,
instead of requiring each application to rediscover how precision grows
through a long iteration.

The long-term integration point is a collection of reusable, proved
finite contracts with explicit demand functions. Each existing packet
can expose which of these contracts it supplies, which it consumes,
and which it still lacks. That organization makes further mathematical
progress concrete: a missing quantitative theorem remains visible,
while previously established finite algebra and analytic estimates
become reusable derivations.


\begin{thebibliography}{9}
\bibitem{ProveIt}
V. Reshetnikov and contributors, \emph{ProveIt},
\path{Logic/PrimitiveRecursiveArithmetic}, inspected at commit
\texttt{27a07921176bccbf630084b959700605783ff4aa}.
\url{https://github.com/VladimirReshetnikov/ProveIt/tree/27a07921176bccbf630084b959700605783ff4aa/Logic/PrimitiveRecursiveArithmetic}.
The accompanying source manifest identifies the files used in this report.

\bibitem{Avigad2002}
J. Avigad, \emph{Weak theories of nonstandard arithmetic and analysis},
in S. G. Simpson (ed.), \emph{Reverse Mathematics 2001},
Lecture Notes in Logic 21, A K Peters, 2005, pp.~19--46.
Preprint dated 2 September 2002.
\url{https://www.andrew.cmu.edu/user/avigad/Papers/nonstandard2.pdf}.

\bibitem{ProofMining}
U. Kohlenbach and P. B. Oliva,
\emph{Proof Mining: A Systematic Way of Analysing Proofs in Mathematics},
BRICS Report RS-02-31, June 2002.
\url{https://www.brics.dk/RS/02/31/BRICS-RS-02-31.pdf}.

\bibitem{Briseid}
E. M. Briseid, \emph{On Rates of Convergence in Metric Fixed Point Theory},
doctoral thesis, Technische Universit\"at Darmstadt, 2009.
\url{https://www2.mathematik.tu-darmstadt.de/~kohlenbach/Briseid-phd.pdf}.
\end{thebibliography}
\end{document}
